% ======================================================================
%  33_flesh_dmod.tex — 骨血层：第 11--12 节
%  主题：典范余各向同性 A-膜作为「翻译器」：膜 → D-模（主链 S8）
%  底本：KW 2006，§11 = 印刷 pp.184--206，§12 = 印刷 pp.207--222
%  锚点约定：source/text/p{N+1:03d}.txt 对应印刷页 N
% ======================================================================

\section{第 11--12 节：典范余各向同性 A-膜作为翻译器}
\label{sec:flesh-dmod}

\begin{keyidea}
\textbf{本章要交付什么。}全文主链到这里还差最后一环。第 8 节已经证明：\(\pi_1(C)\to\LdualC\) 的同态对应 electric eigenbrane，其 S-对偶像是卷在 Hitchin 纤维上的 magnetic eigenbrane（主链 S6）；第 9--10 节又证明 't Hooft 线算子就是几何 Hecke 算子（S7）。可是几何 Langlands 的标准表述要求右边的对象是 \(\MC\) 上的 \term{D-module}{D-模}。本章（S8）说明两件事：第一，第 11 节用复结构 \(K\) 中一个\emph{由整体数据唯一确定}的 A-膜——\term{canonical coisotropic brane}{典范余各向同性膜}（c.c. 膜）——把「\(K\) 中 A-膜」翻译成「\(\MC\) 上的 D-模」；第二，第 12 节用四维规范理论的边界条件说明三类关键膜 \((B,B,B)\)、\((B,A,A)\)、\((A,B,A)\) 确实能在底层规范理论中定义，从而验证整幅图景自洽。对应骨架编号：主链 S8、节点 N11--N12、miracle M5。
\end{keyidea}

\subsection*{本章的符号约定}

除本册固定符号外，本章新出现或改变了用法的符号如下（含义后的括号给出首现锚点）。

\begin{center}
\small
\begin{tabular}{@{}ll@{}}
\toprule
\tabhead{符号} & \tabhead{含义与首现锚点} \\
\midrule
\(\hat N\) & \(\omega^{-1}F\)，c.c. 膜给出的概复结构（式 (11.1)，src p.186）\\
\(I,J,K\) & \(\MH\) 超凯勒结构的三个复结构；\(\omega_I,\omega_J,\omega_K\) 为对应 Kähler 形式（src p.52）\\
\(\Omega_I\) & \(I\) 中全纯辛形式，\(\Omega_I=\omega_J+\ii\omega_K\)，故 \(\Re\Omega_I=\omega_J\)（src p.50, p.189）\\
\(F,\ B\) & Chan-Paton 曲率与 B-field；扭化后进入组合 \(F+B\)（src p.187, p.202）\\
\(\Psi\) & canonical parameter；实 \(\Psi\) 情形取 \(\Psi=\Re\tau=\theta/2\pi\)（src p.202）\\
\(\mathcal N\) & c.c. 膜 Chan-Paton 的平凡线丛 \(\mathcal N\to\MH\)（src p.187）\\
\(L\) & \(\MC\) 上的行列式线丛，其第一 Chern 类为 \(\omega_I/2\pi\)（src p.203, p.206）\\
\(K_M\) & \(M=\MC\) 的典范丛；\(K_M^{1/2}\) 为自旋丛（src p.195）\\
\(\Dmod_L\) & 线丛 \(L\) 上全纯微分算子构成的层（src p.193）\\
\(I(\Psi)\) & 扭化复结构 \(\hat N=\omega^{-1}(F+B)=I\cos q-J\sin q\)（式 (11.32)，src p.204）\\
\(\mathcal B_{c.c.}\) & 典范余各向同性膜；\(\mathcal B_{c.c.}^{\Psi}\) 为其扭化类比（src p.187, p.203）\\
\((\mathrm{A},\mathrm{B},\mathrm{A})\) & 相对于复结构 \((I,J,K)\) 的 A/B 类型三元组（src p.208）\\
\bottomrule
\end{tabular}
\end{center}

\subsection{最后一步的缺口：第 8 节拿到了什么，还缺什么}
\label{subsec:gap}

先给结论。第 8 节证明的是「electric eigenbrane \(\leftrightarrow\) \(\pi_1(C)\to\LdualC\) 的同态」；几何 Langlands 标准表述的右边却是「\(\MC\) 上的 D-模 \(\leftrightarrow\) \(\pi_1(C)\to\LdualC\) 的同态」。两头说的是同一批局部系统，中间却隔着一个范畴：膜是二维拓扑场论的边界条件，D-模是代数几何的对象。原文对此的措辞是一句直白的自陈：

\begin{origquote}
A crucial ingredient is missing from what we have presented so far. \dots\ Hecke eigensheaves on \(M(G,C)\) which are also D-modules.
\end{origquote}
\anchor{src p.185}

\paragraph{缺口有三个层次，缺一不可。}把「差一个环节」拆开，才能看出第 11 节为什么必须做那么多事。

\begin{enumerate}
\item \textbf{对象类型不同。}膜是 A-模型（或 B-模型）的边界条件；D-模是「\term{differential operator}{微分算子}层的模」。前者是物理对象，后者是代数对象。若只说「膜给出 D-模」，必须说明用一个什么样的膜、通过什么运算把边界条件变成层。\pA\anchor{src p.185}
\item \textbf{所在空间不同。}第 8 节的膜生活在 Hitchin 模空间 \(\MH(G,C)\)；D-模生活在 \(\MC=M(G,C)\)，即 \(C\) 上稳定 \(G\)-丛的模空间。两者维数相同（都是 \(2\dim M(G,C)\) 与 \(\dim M(G,C)\) 的关系），但不是同一个空间：\(\MH\) 是 \(\MC\) 的余切丛（在稳定点附近）。\pA\anchor{src p.53}
\item \textbf{类型不匹配。}由第 8--10 节得到的 magnetic eigenbrane 是卷在 Hitchin 纤维上的膜，它是 \((A,B,A)\)-膜；而 \((A,B,A)\) 中的「A」只相对于 \(I\) 与 \(K\)，「B」只相对于 \(J\)。标准 Langlands 表述并不提复结构 \(I,J,K\)，所以还要说明：在哪一个 A-模型里做翻译，翻译出来的东西为什么就是标准的 D-模。\pA\anchor{src p.185}
\end{enumerate}

\paragraph{原文给的路线：把 D-模看成非交换形变上的 B-膜。}原文并不打算从零发明一个「膜到 D-模」的函子；它引用了一个更早的观察\pA\anchor{src p.185}：\(\MC\) 上的 D-模可以看作 \(T^*\MC\) 的\emph{非交换形变}上的 B-膜；而对某一类全纯辛流形，A-膜范畴与「同一流形非交换形变上的 B-膜范畴」等价（[41]）。于是本节的翻译可以看作那个等价的特例。我们判断，这个「特例」的说法是全文归约路线的核心：正因为它是一个特例，第 11 节才只需要构造一个具体的膜，而不需要一般地发展非交换几何。\pB\anchor{src p.185}

\caveat{原文紧接着补了一句常被忽略的话：在几何 Langlands 中只做在稳定丛模空间 \(\MC\) 上是不够的，Hecke 算子的作用方式要求把 D-模做到 \(\mathrm{Bun}_G(C)\)（所有 \(G\)-丛的 stack）上；规范理论里的类比是「所有联络的空间」，这正是第 12 节要说明的。\pA\anchor{src p.185}}

\paragraph{本章的推进顺序。}图~\ref{fig:flow} 把第 11--12 节的逻辑一次画完：第 11.1 节先造出 \(\mathcal B_{c.c.}\)；第 11.2 节算出它的开弦代数并证明它逐点同构于全纯微分算子层，于是「\(K\) 中 A-膜」自动成为该代数层的模，翻译完成；第 11.3 节把 \(\Psi=0\) 推广到任意 \(\Psi\)，得到扭化微分算子；第 12 节回到四维规范理论，用边界条件重新造出三类关键膜，作为整幅图景的自洽性检查。

\begin{figure}[htbp]
\centering
\begin{tikzpicture}[font=\small,
  box/.style={draw=kwblue!70, rounded corners=2pt, fill=kwboxbg, align=center,
              inner sep=4pt, text width=4.6cm, minimum height=1.1cm},
  src/.style={draw=kwgray!70, rounded corners=2pt, fill=white, align=center,
              inner sep=3pt, text width=3.6cm},
  arr/.style={-{Stealth[length=2.2mm]}, thick, draw=kwblue!70}]
  \node[src] (s6) at (-5.2,3.4) {S6（\S8）：electric eigenbrane\\ \(\leftrightarrow\) \(\pi_1(C)\to\LdualC\)};
  \node[src] (s7) at (5.2,3.4) {S7（\S9--10）：'t Hooft 算子\\ \(\leftrightarrow\) Hecke 算子};
  \node[box] (cc) at (0,1.9) {S8 第 1 步（\S11.1）：造出\\ 典范余各向同性 A-膜 \(\mathcal B_{c.c.}\)};
  \node[box] (alg) at (0,0.35) {S8 第 2 步（\S11.2）：开弦代数\\ \(\cong\Dmod_{K_M^{1/2}}\)，于是每个\\ \(K\) 中 A-膜 \(\mapsto\) D-模};
  \node[box] (tw) at (0,-1.2) {S8 第 3 步（\S11.3）：扭化\\ \(\Psi\neq0\) 时 \(\Dmod_{K_M^{1/2}\otimes L^{\Psi}}\)};
  \node[box] (gl) at (0,-2.75) {标准几何 Langlands：\\ \(\MC\) 上 D-模 \(\leftrightarrow\pi_1(C)\to\LdualC\)};
  \node[src] (s12) at (0,-4.3) {\S12 检查清单：四维边界条件\\ \((B,B,B)\)、\((B,A,A)\)、\((A,B,A)\)};
  \draw[arr] (s6) -- (cc);
  \draw[arr] (s7) -- (cc);
  \draw[arr] (cc) -- (alg);
  \draw[arr] (alg) -- (tw);
  \draw[arr] (tw) -- (gl);
  \draw[arr, dashed] (s12) -- (gl);
\end{tikzpicture}
\caption{第 11--12 节在全文中承担的论证任务。实线是「导出」关系，虚线是「自洽性检查」关系：第 12 节不产生新结论，它验证第 11 节造出的三类膜确实可以在四维规范理论中定义。}
\label{fig:flow}
\end{figure}

\paragraph{排除错误直觉 1：D-模不是「把膜量子化」得到的产物。}常见的错误图像是：仿佛存在一个叫「量子化」的函子，把膜范畴变成 D-模范畴。原文的路线恰恰相反。它先在 \(K\) 的 A-模型里算出一个具体膜 \(\mathcal B_{c.c.}\) 的\emph{开弦代数}（(B,B) 弦构成的代数），发现这个代数逐点同构于全纯微分算子层；于是「D-模」不是从膜范畴新造出来的结构，而是同一批开弦态的另一种说法：代数层的模就是别的膜上的开弦态。翻译是「识别」而非「量子化」。\pA\anchor{src p.188}\pC\anchor{src p.197}

\subsection{典范余各向同性 A-膜：存在、唯一与「典范」的含义（\S11.1）}
\label{subsec:cc}

\paragraph{动机：翻译器为什么必须「空间充满」。}设我们要把任意一张膜 \(B'\) 翻译成一个模。按照开闭拓扑场论的公理，对任意两张膜 \(B,B'\)，\((B,B')\) 弦构成 \((B,B)\) 弦代数上的一个模；而 \((B,B)\) 弦就是 \(B\) 作为膜范畴对象的 endomorphism 代数。\pA\anchor{src p.188} 因此，只要找到一张\emph{空间充满}的膜 \(B\)，并把它的 endomorphism 代数算清楚，所有膜就自动成为该代数的模。麻烦在于：普通的 A-膜支撑在 Lagrangian 子流形上，维数是 \(\dim X/2\)，不可能是空间充满的（除非 \(\dim X=0\)）。所以翻译器的第一要求是：它必须是 A-膜，但支撑维数高于中维。

这正是文献 [80] 引入的\term{coisotropic A-brane}{余各向同性 A-膜}：支撑 \(Y\) 由局部一堆 Poisson 交换的约束函数 \(f_1,\dots,f_d\) 的零点定义（第一类约束），故 \(\dim Y\ge \dim X/2\)；取等号时 \(Y\) 是 Lagrangian，回到普通 A-膜；高于中维的情形一般存在，但条件特殊。\pA\anchor{src p.186} 第 11 节只需要最简单的情形：\(Y=X\)、Chan-Paton 丛是 rank one 的 \(U(1)\) 线丛。

\paragraph{派生对象：\(U(1)\) 规范场与条件 (11.1)。}取 Chan-Paton 线丛 \(L\to X\)；\(U(1)\) 规范场局部是实 1-形式 \(A\)，协变导数 \(D=\dd+A\)，曲率 \(F=\dd A\)；取 \(U(1)\) 的 \(n=1\) 表示，世界面上与 Chan-Paton 丛相伴的因子是边界 holonomy 的指数 \(\exp(\ii\oint_{\partial\Sigma}\Phi^*A)\)。\pA\anchor{src p.186} 原文断言：这样的膜要是 A-膜，其曲率必须满足
\begin{equation}
(\omega^{-1}F)^2=-1,
\label{eq:cc-def}
\end{equation}
即 \((1,1)\)-张量 \(\hat N=\omega^{-1}F\) 是一个概复结构；又因 \(F,\omega\) 闭，(11.1) 与它们合起来蕴含 \(\hat N\) 可积。\pA\anchor{src p.186}

原文在这里只引用 [80]，没有给出推导。下面是我们补的重建（\pB），它把「为什么偏偏是 \(\hat N\)」讲清楚；重建的依据是原文紧接着用的三条性质，它们都可以由 (11.1) 与 \(\dd\omega=\dd F=0\) 直接算出。

\paragraph{推导 A：条件 (11.1) 的几何内容。}记 \(\hat N=\omega^{-1}F\)，等价地
\begin{equation}
F(X,Y)=\omega(\hat N X,Y)\qquad \text{对任意向量场 }X,Y.
\label{eq:F-N}
\end{equation}
若 \(\hat N^2=-1\)，则 \(\hat N\) 给出切丛的概复结构。我们先算 \(\omega\) 与 \(F\) 相对于 \(\hat N\) 的类型：对任意 \(X,Y\)，
\begin{align}
F(\hat N X,Y) &= \omega(\hat N^2 X,Y)=-\omega(X,Y), \label{eq:type-1}\\
F(\hat N X,\hat N Y) &= \omega(\hat N^2X,\hat N Y)=-\omega(X,\hat N Y)=-F(X,Y). \label{eq:type-2}
\end{align}
由 \eqref{eq:type-1} 与 \eqref{eq:type-2} 知 \(\omega\) 与 \(F\) 都是 \((2,0)\oplus(0,2)\) 型。再把它们按 \((p,q)\) 分解：在 \((1,0)\) 向量上 \(\hat N=\ii\)，在 \((0,1)\) 向量上 \(\hat N=-\ii\)。于是对两个 \((1,0)\) 向量，
\begin{equation}
F^{2,0}(X,Y)=F(X,Y)=\omega(\hat N X,Y)=\ii\,\omega(X,Y)=\ii\,\omega^{2,0}(X,Y),
\label{eq:type-3}
\end{equation}
故 \(\omega-\ii F\) 的 \((2,0)\) 部分为 \(2\omega^{2,0}\)；类似地，对两个 \((0,1)\) 向量有 \(F^{0,2}=-\ii\omega^{0,2}\)，从而 \(\omega-\ii F\) 的 \((0,2)\) 部分为零。因此
\begin{equation}
\omega-\ii F\ \text{是} (2,0)\text{型，且为全纯辛形式。}
\label{eq:hol-symp}
\end{equation}
反之，「\(\omega-\ii F\) 为全纯辛形式」也强迫 \(\hat N^2=-1\)，所以 (11.1) 不是外加条件，而是「\(\omega\) 的复化与 \(F\) 一起给出全纯辛结构」的等价说法。\pB\anchor{src p.186}

这条推导立刻给出原文的一个推论：\emph{这种余各向同性 A-膜只可能在 \(\omega\) 是某个全纯辛形式的实部时才存在}。\pA\anchor{src p.186} 理由很直接：由 \eqref{eq:hol-symp}，\(\omega=\Re(\omega-\ii F)\)；若 \(X\) 上不存在使 \(\omega\) 成为全纯辛形式实部的复结构，条件 (11.1) 无解。

\paragraph{在 \(\MH\) 上求解 (11.1)。}取 \(X=\MH(G,C)\)，A-模型在复结构 \(K\) 中，辛形式是
\begin{equation}
\omega=(\Im\tau)\,\omega_K .
\label{eq:omega-K}
\end{equation}
把 (11.1) 中取 \(\hat N=\omega^{-1}F\)，显然有解
\begin{equation}
F=(\Im\tau)\,\omega_J,\qquad \hat N=\omega_K^{-1}\omega_J=I,
\label{eq:F-sol}
\end{equation}
其中 \(\hat N=I\) 用到四元数关系 \(IJ=K\)、\(JK=I\)、\(KI=J\)（以及 \(\omega_K^{-1}\omega_J=K^{-1}J=I\)）。另有一个 \(F\) 与 \(\hat N\) 同时变号的解。\pA\anchor{src p.187}

更一般地，原文考虑
\begin{equation}
F=(\Im\tau)\,(a\,\omega_J+b\,\omega_I),\qquad a^2+b^2=1,
\label{eq:family}
\end{equation}
它同样解 (11.1)：直接算得
\begin{equation}
\hat N=\omega_K^{-1}F=aI+b\,\omega_K^{-1}\omega_I=aI-bJ,
\qquad \hat N^2=-(a^2+b^2)\,\mathbf 1.
\label{eq:N-family}
\end{equation}
但原文指出：\(b\neq0\) 时这个形式\emph{不恰当}，而且对一般 \(\tau\) 其周期不是 \(2\pi\) 的整数倍，因此通常不是任何联络的曲率，不能用来造 A-膜。\pA\anchor{src p.187}

我们把这句话的算术补全（\pB）。两件事分开看。第一，\(F\) 必须是平凡线丛 \(\mathcal N\) 上某个联络的曲率，即 \(F=\dd\alpha\) 整体恰当；这一点在 \(b=0\) 时成立（\(\omega_J\) 在 \S4.1 已证恰当，见原文同页），在 \(b\neq0\) 时不成立：\(\omega_I\) 有非零周期——它限制到 Hitchin 纤维（复环面）上仍是非退化的闭 2-形式，而环面的 \(H^2\) 非零。第二，即使只要求 \(F\) 是某个 \(U(1)\) 线丛的曲率，也需要周期落在 \(2\pi\ZZ\) 中；对于一般的 \(\tau\)，\((\Im\tau)b\oint\omega_I\) 显然不满足。于是 (11.4) 中去掉 \(b=0\) 的两个解后，其余都不是合法的 A-膜。\pB\anchor{src p.187}

\begin{definition}[典范余各向同性膜]
\label{def:cc-brane}
在 \(\MH(G,C)\) 上取 A-模型、辛形式 \(\omega=(\Im\tau)\omega_K\)，Chan-Paton 丛为平凡线丛 \(\mathcal N\to\MH\)。令 \(F=(\Im\tau)\omega_J\)，它是 \(\mathcal N\) 上联络的曲率形式；满足 (11.1) 的空间充满膜称为\term{canonical coisotropic brane}{典范余各向同性膜}，简称 c.c. 膜，记作 \(\mathcal B_{c.c.}\)。
\end{definition}

原文紧接着给的定义与唯一性断言是：「\(\MH\) 上 \(F=(\Im\tau)\omega_J\) 的 A-膜称为 c.c. 膜；若 \(G\) 单连通则 \(\MH\) 也单连通，c.c. 膜唯一；否则唯一到把 Chan-Paton 丛用平坦线丛扭一下。」\pA\anchor{src p.187} 注意这里的措辞：唯一性是在「\(F\) 固定」之后说的；确定 \(F\) 的是 \(\omega_J\)，而 \(\omega_J=\Re\Omega_I\) 是由 \(\MH\) 的复结构 \(I\) 整体决定的，不涉及任何额外选择。这正是「典范」二字的含义。

\paragraph{\(\mathcal B_{c.c.}\) 是 \((A,B,A)\)-膜。}原文的第一个观察是：c.c. 膜在复结构 \(I\) 与 \(K\) 中是 A-膜，在复结构 \(J\) 中是 B-膜。\pA\anchor{src p.187} 我们把它拆成三步验证。

\begin{itemize}
\item \textbf{在 \(I\) 中是 A-膜。}只需验证 \(\hat N_I=\omega_I^{-1}F=(\Im\tau)\omega_I^{-1}\omega_J\) 满足 (11.1)。用四元数关系 \(\omega_I^{-1}\omega_J=-K\)，故
\begin{equation}
(\omega_I^{-1}\omega_J)^2=(-K)^2=-1 .
\label{eq:ABCA-I}
\end{equation}
这就是原文的 (11.5)；\(F\) 是实 2-形式且闭，由推导 A 的同一论证，\(\hat N_I\) 可积，因此它确实满足 \(I\) 中余各向同性 A-膜的条件。\pA\anchor{src p.187}
\item \textbf{在 \(K\) 中是 A-膜。}这正是 \eqref{eq:F-sol} 的构造。
\item \textbf{在 \(J\) 中是 B-膜。}B-膜的两个条件是：支撑在 \(J\) 中全纯（\(X\) 自身当然是）、Chan-Paton 丛在 \(J\) 中全纯。后者等价于曲率的 \((0,2)\) 部分为零。现在 \(F=(\Im\tau)\omega_J\) 是 \(\omega_J\) 的倍数，而 \(\omega_J\) 在 \(J\) 中是 \((1,1)\) 型，故 \(F^{0,2}=0\)：由 \(F=\bar\partial_A^2\)（在 \(J\) 中）知 \(\bar\partial_A\) 定义了一个全纯结构。\pA\anchor{src p.187}
\end{itemize}

顺带一个原文没有写、但由 \eqref{eq:N-family} 立即得到的小结论，它解释了 \((A,B,A)\) 这个标签为什么是「被逼出来的」而不是「选出来的」：在 (11.4) 这一自然族里，若要求同时在 \(K\) 与 \(I\) 中为 A-膜，则 \(a^2+b^2=1\) 与 \(b=0\) 联立给出 \(a=\pm1\)，此时 \(\hat N_J=a+bK=\pm1\)，其平方为 \(+1\) 而不是 \(-1\)，故 \(J\) 中必为 B-膜。也就是说，\((A,A,A)\) 在这族里不可能出现。\pC\anchor{eq 11.4}

\caveat{公式 (11.2) 只在 \(\MH\) 的光滑部分定义 c.c. 膜；原文明确说，它可以在底层规范理论中由一个边界条件定义（见 \S12.4），并强调这句陈述正是「构造在整个 \(G\)-丛 stack 上成立」的量子场论等价说法。换句话说，c.c. 膜不是 \(\MH\) 的 sigma 模型里一个碰巧写出来的对象，它有四维来源。\pA\anchor{src p.187}}

\paragraph{排除错误直觉 2：c.c. 膜不是「随便挑一个余各向同性膜」。}读原文容易产生一种错觉：既然 (11.4) 有一整族解，c.c. 膜只是其中参数取 \((a,b)=(1,0)\) 的一个特例，那么「典范」二字从何而来？错在三处。第一，(11.4) 的其余元素大多数根本不是合法的膜——它们不恰当、周期不量子化，连 Chan-Paton 线丛都造不出来；所以可选的解集只有两个点：\((1,0)\) 与 \((-1,0)\)，而后者只是整体变号。第二，参数不只是「选」出来的：\(\omega_J=\Re\Omega_I\) 是复结构 \(I\) 附带的全纯辛形式，它的实部是唯一确定的，因此 \(F\) 唯一确定。第三，如果 \(G\) 不是单连通，唯一性要打折扣，但折扣的方式也是典范的：只能用\emph{平坦}线丛去扭，而平坦扭不改变开弦代数层（见 \S11.2 中关于 \(\Dmod_L\) 的讨论）。所以「典范」= 「由 \(\MH\) 的整体复几何唯一决定，且与所有任意选择无关」。\pA\anchor{src p.187}\pC\anchor{src p.193}

\paragraph{排除错误直觉 3：\(F\neq0\) 不表示「膜上有非平凡的动力学规范场」。}另一种误读是把 c.c. 膜的 \(F=(\Im\tau)\omega_J\) 当作「给膜加了一个强背景规范场，所以它和普通膜很不一样」。事实相反：\(\mathcal N\) 是\emph{拓扑平凡}的线丛，\(F\) 只是这个平凡丛上一个固定的联络，由 \(\MH\) 的几何决定，不是可变的动力学自由度。真正关键的是：这个固定的 \(F\) 使 \(\hat N=\omega^{-1}F\) 成为复结构 \(I\)，从而改变了「开弦代数里的复结构」；\(F\) 的拓扑平凡性恰好保证了这种改变不引入新的拓扑选择。\pA\anchor{src p.187}

\begin{enumerate}
\item 读到这里的自检：如果 \(\omega\) 不是全纯辛形式的实部，条件 (11.1) 为什么不可能有解？请用 \(\omega-\ii F\) 的类型论证。
\item c.c. 膜与「\(F=0\) 的空间充满膜」在 \(K\) 中各是什么类型的膜？为什么后者不能当翻译器？
\end{enumerate}

\subsection{开弦代数就是 D-模：翻译机制（\S11.2）}
\label{subsec:translation}

\paragraph{动机：不需要发明一个函子，只需要算一个代数。}上一小节造出了空间充满的 \(\mathcal B_{c.c.}\)。按开闭拓扑场论的公理，对任意另一张膜 \(B'\)，\((B',B_{c.c.})\) 弦构成 \((B_{c.c.},B_{c.c.})\) 弦代数上的模；所以只要把 \((B_{c.c.},B_{c.c.})\) 弦的代数算出来，翻译就自动发生。\pA\anchor{src p.188} 第 11.2 节的全部技术工作因此只服务于一个目标：证明这个代数逐点同构于全纯微分算子层 \(\Dmod_{K_M^{1/2}}\)。原文的表述是：\((B,B)\) 弦是膜 \(B\) 作为膜范畴对象的 endomorphism 代数；我们要在 \(B=\mathcal B_{c.c.}\) 的情形把它算出来。\pA\anchor{src p.188}

\paragraph{第一步：世界面瞬子没有贡献，所以边界可观测量可以层化。}计算 \((B_{c.c.},B_{c.c.})\) 弦之前，必须先排除非微扰贡献。原文先写下 [80] 给出的余各向同性膜边界条件（\(\partial_1,\partial_0\) 分别是世界面边界的法向与切向导数）：
\begin{equation}
G^{IJ}\partial_1\Phi_J=-\ii F^{IJ}\partial_0\Phi_J,
\label{eq:coiso-bc}
\end{equation}
再要求世界面映射在复结构 \(K\) 中全纯（即瞬子方程），把 \(\partial_1\Phi\) 用 \(\partial_0\Phi\) 表出，得
\begin{equation}
(\omega-\ii F)\partial_0\Phi=0.
\label{eq:instanton-kill}
\end{equation}
由于 \(\omega\) 与 \(F\) 都非退化，\(\partial_0\Phi=0\)：瞬子沿边界为常值，再由解析性处处为常值，故是平凡瞬子。\pA\anchor{src p.188}

没有瞬子修正，世界面微扰论逐阶产生的修正都在 target 空间局部；于是「边界可观测量」的合适对象不是单个代数，而是 target 上的层。\pA\anchor{src p.188} 在经典极限，任意余各向同性 A-膜 \(Y\) 的 BRST 不变边界算子空间已由 [80,132] 确定；只取 ghost number 0，在 \(Y=X\) 的情形它们就是复结构 \(\hat N\) 下的全纯函数。\pA\anchor{src p.188} 对 c.c. 膜，\(\hat N=I\)，所以可加地（先忘掉环结构）：
\begin{keyidea}
c.c. 膜的拓扑边界算子层 \(=\MH\) 上对复结构 \(I\) 全纯的函数层。
\end{keyidea}
\noindent 原文特别提醒：我们虽然是在复结构 \(K\) 的 A-模型里工作，真正相关的复结构却是 \(\hat N=I\)。\pA\anchor{src p.189}

\paragraph{第二步：从对易代数到非交换代数。}经典层是交换的；量子修正来自弦的延展性。原文用非交换参数 \(\theta\)（文献 [137] 的记号）说明这一点：\(\theta\) 是 \((g+F)^{-1}\) 的反对称部分，在本问题中成为 \(\omega_J^{-1}\) 的倍数；而 \(\omega_J=\Re\Omega_I\)，当我们作用在「对 \(I\) 全纯的函数」（即拓扑可观测量）上时，可以把 \(\omega_J^{-1}\) 换成 \(\Omega_I^{-1}\)。于是乘法的形变由 \(\Omega_I\) 给出的 Poisson 括号控制。\pA\anchor{src p.189}

这里有一个物理上容易被误读的细节，原文交代得很清楚：在通常的物理弦论中，只有当 \(|F|\gg|g|\) 时开弦代数才有一个好的非交换形变描述；本题中 \(|F|=|g|\)，并不在该极限内。然而限制到拓扑可观测量之后会发生与 \(|F|\gg|g|\) 时类似的简化——在 B-型拓扑弦中这一点首见于 [71]；物理上靠 \(|F|\gg|g|\) 消去的弦激发，在这里被「只取拓扑可观测量」消去了。\pA\anchor{src p.189}

拓扑理论在圆盘上的作用量可写成 \(S=\int_\Sigma\Phi^*(\omega-\ii F)+\int_\Sigma\dd^2\sigma\{Q,V\}\)（11.8）。由于
\begin{equation}
\omega-\ii F=(\Im\tau)\omega_K-\ii(\Im\tau)\omega_J=-\ii(\Im\tau)\Omega_I,\qquad \Omega_I=\dd\varpi,
\label{eq:S-theta}
\end{equation}
其中 \(\varpi\) 是 \(\MH\) 上典范全纯 1-形式 (11.10)，所以模掉 \(\{Q,V\}\) 后作用量是边界项 \(S=-\ii\Im\tau\oint_{\partial\Sigma}\Phi^*\varpi\)（11.9）。\pA\anchor{src p.189}\pA\anchor{src p.190} 限制到 \(T^*\MC\)：取 \(q^\alpha\) 为 \(\MC\) 局部全纯坐标、\(p_\alpha\) 为纤维线性函数，则（差数值因子）\(\varpi=\sum_\alpha p_\alpha\dd q^\alpha\)，作用量成为 \(S=-\ii\Im\tau\int p_\alpha\dd q^\alpha+\{Q,\dots\}\)（11.10 后），即「\(\MC\) 上零 Hamilton 量粒子量子力学」的全纯版本。\pA\anchor{src p.190} 这个限制丢掉 \(\MH\) 上没有的一般结构，但保留下需要的全部结构：全纯辛流形、C$^*$ 作用、沿纤维的量子化。

\paragraph{推导 B（完整）：从圆盘作用量到 \(\Dmod_{K_M^{1/2}}\)。}下面把 (11.11)--(11.16) 之间被原文压缩的步骤补全。这条推导较长，分五步；每一步我们都注明是原文明确写出（\pA）还是我们补的（\pB）。

\emph{第 1 步：C$^*$ 分级与局部性。}原文用 \S4.2 末的 C$^*$ 作用 \(U(\lambda):\varphi_z\mapsto\lambda\varphi_z\)（式 (11.11)）给可观测量分次：势 \(\varpi\) 的 degree 为 1，若把 \(1/\Im\tau\)（相当于 Planck 常数）也赋 degree 1，则 \(U\) 提升为量子对称。\pA\anchor{eq 11.11}\pA\anchor{src p.190} 于是只在 C$^*$-不变开集 \(V\subset\MC\) 上、并沿 \(T^*V\) 的纤维要求多项式增长来研究代数：局部代数由 degree \(\le1\) 的元素生成，低次量子修正只让代数成为 filtered 而非 graded。\pA\anchor{src p.191}

\emph{第 2 步：degree 0 与 degree 1 元素。}degree 0 是 \(V\) 上全纯函数 \(f\)（没有负 degree 可观测量，故不受形变影响）；degree 1 是沿纤维线性的函数 \(\kappa=\sum_\alpha p_\alpha V^\alpha(q)\)，对应全纯向量场 \(\sum_\alpha V^\alpha\partial_{q^\alpha}\)。两者的乘积只有一阶形变，原文的 (11.12) 是
\begin{equation}
[\hat\kappa,\hat f]=-\frac{\ii}{\Im\tau}\,\widehat{\kappa(f)},\qquad \kappa(f)=\sum_\alpha V^\alpha\frac{\partial f}{\partial q^\alpha}.
\label{eq:comm-kf}
\end{equation}
\pA\anchor{src p.192} 我们用正则量子化把它再算一遍（\pB）：设 \([\hat p_\alpha,\hat q^\beta]=-(\ii/\Im\tau)\delta_\alpha^{\ \beta}\)，则在 \(\hat\kappa=\sum_\alpha V^\alpha(\hat q)\hat p_\alpha\) 的 Weyl 排序下，(11.12) 等价于
\begin{equation}
[\hat p_\alpha,\hat f]=-\frac{\ii}{\Im\tau}\frac{\partial f}{\partial q^\alpha},
\qquad\text{即}\qquad
\hat p_\alpha=-\frac{\ii}{\Im\tau}\frac{\partial}{\partial q^\alpha}
\quad\text{（在 }T^*V\text{ 上）},
\label{eq:p-hat}
\end{equation}
这正是原文的 (11.15)。\pA\anchor{src p.192}

\emph{第 3 步：冗余与二阶形变。}只保留 degree \(\le1\) 的元素时存在冗余 \(\hat p_\alpha\mapsto\hat p_\alpha+A_\alpha(q)\)（(11.14)），它对应不改变 \(\varpi\) 的高阶项。\pA\anchor{src p.192} 两个 degree 1 元素相乘时一阶形变仍由 Poisson 括号给出，但还可能留下二阶形变 \(\sim\hbar^2F_{\alpha\beta}(q)\)；由结合律可证 \(F_{\alpha\beta}\dd q^\alpha\dd q^\beta\) 闭，在小开集 \(V\) 上恰当，用 (11.14) 的冗余可化零。故 \(T^*V\) 上的代数就是 \(V\) 上全纯微分算子代数。\pB\anchor{src p.192}

\emph{第 4 步：粘起来，得到 \(\Dmod_L\) 或 \(\Dmod_{L^\gamma}\)。}用 C$^*$-不变开集 \(V^{(i)}\) 覆盖 \(M=\MC\)：重叠区 degree 0 元素必须一致，degree 1 元素可差 (11.14) 型变换；再由每个开集内 \([\hat p_\alpha,\hat p_\beta]=0\)，必须有
\begin{equation}
\hat p_\alpha^{(i)}=\hat p_\alpha^{(j)}
+\frac{1}{\Im\tau}\frac{\partial s^{(ij)}}{\partial q^\alpha}
\qquad\text{在 }V^{(i)}\cap V^{(j)}\text{ 上},
\label{eq:glue}
\end{equation}
其中 \(s^{(ij)}\) 全纯。\pA\anchor{src p.193} 三重交上的量 \(c^{(ijk)}=s^{(ij)}+s^{(jk)}+s^{(ki)}\) 是复常数、自动满足 cocycle 条件，定义 \(H^2(M,\CC)\) 中元素：若落在 \(2\pi\ZZ\) 中，则 \(\exp(-\ii s^{(ij)})\) 是复线丛 \(L\to\MH\) 的转移函数，开弦代数层就是 \(\Dmod_L\)；若 \(L\) 拓扑非平凡而常数只是任意复数，则不能乘 \(\gamma\) 定义新线丛，但可以把 \(s^{(ij)}\) 乘 \(\gamma\) 定义层 \(\Dmod_{L^\gamma}\)。\pA\anchor{src p.193} 原文指出 \(H^2(\MC,\CC)\) 一维，故本问题最一般的层是 \(\Dmod_{L^\gamma}\)（\(\gamma\in\CC\)，\(L\) 为行列式线丛）。\pA\anchor{src p.193} 此分析无微扰收敛问题（二阶决定）；完全层化才有收敛困难（C$^2$ 上 Moyal 乘积的例子）。\pA\anchor{src p.191}\pA\anchor{src p.193}

\emph{第 5 步：时间反演把 \(L\) 钉死为 \(K_M^{1/2}\)。}到这一步为止，代数层是某个 \(\Dmod_L\)，但 \(L\) 是哪一个还没定；原文用时间反演 \(T\) 确定它。三件事：（i）\(T\) 恰在 \(t=1\) 是扭化 TFT 的对称（式 (3.56)），而我们正工作在 \(t=1\)；（ii）二维理论中 \(T\) 是世界面 \(\Sigma\) 的反定向对称，把边界映到自身并反转边界定向，玻色子变换为 \((A,\varphi)\mapsto(T^*A,-T^*\varphi)\)（式 (3.57)），\(\varphi\) 前的负号使 (11.8)、(11.9) 中的 \(\omega\)、\(F\)、\(\varpi\) 项不变；（iii）\(T\) 反转算子的插入顺序，故把开弦代数映到 opposite 代数。\pA\anchor{src p.194}

opposite 代数是哪一层？对线丛 \(L\) 上的全纯微分算子代数，分部积分给出转置 \(F\mapsto F^{\mathrm t}\)（\((FG)^{\mathrm t}=G^{\mathrm t}F^{\mathrm t}\)），它把 \(\Dmod_L\) 送到 \(\Dmod_{L^{-1}\otimes K_X}\)；因为 \(F^{\mathrm t}\) 有全纯系数且与实切片 \(Z\) 的选取无关，所以
\begin{equation}
(\Dmod_L)^{\mathrm{op}}\cong\Dmod_{L^{-1}\otimes K_X}.
\label{eq:opposite-dmod}
\end{equation}
\pA\anchor{src p.195} c.c. 膜具有时间反演对称，故其开弦代数同构于自己的 opposite；又已知它形如 \(\Dmod_L\)，于是 \(L\cong L^{-1}\otimes K_X\)，即 \(L\cong K_X^{1/2}\)（\(X=M=\MC\)）。所以
\begin{equation}
\boxed{\ \text{c.c. 膜的开弦代数层}\ =\ \Dmod_{K_M^{1/2}}\ }.
\label{eq:main-ring}
\end{equation}
\pA\anchor{src p.195} 原文补充：\(M\) 实际上是 spin 流形，故 \(K_M^{1/2}\) 是普通线丛；spin 结构不唯一，但「线丛上的微分算子层」只差平坦线丛的扭，故 \(\Dmod_{K_M^{1/2}}\) 不依赖 \(K_M^{1/2}\) 的选取。\pA\anchor{src p.195}

\paragraph{第 6 步（整体截面）：为什么这确实是一个「D-模的环」。}局部代数层还不够，还要看整体截面是否给出正确的对象。原文的答案：整体可观测量 = 层在 \(M\) 上的整体截面（更一般是取值于该层的上同调），只取 ghost number 0、但保留 C$^*$ 的任意标度 degree。\pA\anchor{src p.195} 经典极限下这是 \(\MH\) 上整体全纯函数，正是 \S4.3 中可积系统的交换 Hamilton 量；量子层面则有来自 (11.16) 的胶合修正。\pA\anchor{src p.195} 具体地，给定 \(\MH\) 上一个三次全纯函数 \(w=\sum f^{\alpha\beta\gamma}(q)p_\alpha p_\beta p_\gamma\)，要提升为整体 \((B_{c.c.},B_{c.c.})\) 弦，就要找 \(M\) 上三阶微分算子 \(\hat w\)，其 symbol 为 \(w\)：
\begin{equation}
\hat w=(-\ii/\Im\tau)^3 f^{\alpha\beta\gamma}\frac{\partial^3}{\partial q^\alpha\partial q^\beta\partial q^\gamma}+\cdots,
\label{eq:quantize-w}
\end{equation}
省略号是二阶及以下的项。\pA\anchor{src p.196} 这种量子化存在：degree 2 是 Hitchin 的经典结果 [141]，任意 degree 由 Beilinson--Drinfeld 给出 [5]；原文另给一个用 \(T\) 的论证（见 \S12.4 的算符分类）。\pA\anchor{src p.196} 因此「\(\Dmod_{K_M^{1/2}}\)」不是形式记号：它的整体截面确实把 \(\MH\) 的经典可积系统量子化为 \(M\) 上的量子可积系统。\pA\anchor{src p.196}

\paragraph{第四步：翻译命题。}现在可以把「翻译」写成一个精确的命题。它只说一件事：任何 \(K\) 中的 A-膜都给出一个 D-模。原文在 \S11.2 开头就是这么陈述的（「\(K\) 中每一个 A-膜自动给出 \(\MC\) 上的 D-模」）。\pA\anchor{src p.185}

\begin{proposition}[膜到 D-模的翻译]
\label{prop:translation}
设 \(B'\) 是 \(\MH(G,C)\) 上复结构 \(K\) 中的一张 A-膜。则 \((B_{c.c.},B')\) 弦构成的层是 \(M=\MC\) 上 \(\Dmod_{K_M^{1/2}}\) 的一个模层；等价地，\(B'\) 给出 \(M\) 上一个扭化 D-模。
\end{proposition}

\begin{proof}[证明要点]
（i）\emph{模结构来自公理。}开闭拓扑场论中，任意两张膜 \(B,B'\) 的 \((B,B')\) 弦构成 \((B,B)\) 弦代数的模（图 6 of \S6.3；公理化见 [142,143]）。这是拓扑场论的标准性质，不依赖具体模型。\pA\anchor{src p.197}

（ii）\emph{瞬子仍不贡献。}把 (11.6)--(11.7) 的论证用到 \((B_{c.c.},B')\) 弦：圆盘边界的一段落在 c.c. 膜上，瞬子在这一段上为常值，由解析性处处为常值。故高阶拓扑耦合（例如三次 Yukawa 耦合）也没有非微扰修正。\pA\anchor{src p.197}

（iii）\emph{层化只在 \(T^*V\) 上做。}微扰论局部，因此可以在 \(T^*V\subset\MH\)（\(V\subset M\) 开）上把 \((B_{c.c.},B')\) 弦定义成一个层；沿纤维的多项式增长条件使这个层有好的代数结构。\pA\anchor{src p.197}

（iv）\emph{环层已算出。}由推导 B，\((B_{c.c.},B_{c.c.})\) 弦的环层是 \(\Dmod_{K_M^{1/2}}\)；于是 (iii) 中的层是它的模。\pA\anchor{src p.197}
\end{proof}

\caveat{命题 \ref{prop:translation} 对 \(B'\) 有一个不显眼但必要的前提：除了在 \(K\) 中是 A-膜，它的支撑还必须对 \(\hat N=I\) 复。否则 BRST 算子 \(Q\) 不会约化为 \(I\) 中的 \(\bar\partial\)，\((B_{c.c.},B')\) 弦就不会是微分算子层的模。原文在讨论第一个例子时明确点了这个条件。\pA\anchor{src p.198}}

\paragraph{为什么「余各向同性」恰好给出「可积」的 D-模结构。}把上面的机制抽象一下：c.c. 膜给出一对 \((\omega,F)\)，使 \(\hat N=\omega^{-1}F\) 成为复结构；\(\hat N\) 可积，所以 \(\MH\) 在 \(\hat N\) 下有全纯坐标；对这些坐标，拓扑超荷 \(Q\) 正好是 \(\bar\partial\)。于是「开弦代数的模」自动是「\(\bar\partial\) 的模」，也就是 D-模。反之，如果 \(\hat N\) 不可积，\(Q\) 只是一个形变量子化代数上的算子，我们得到的是非交换代数层，而不是真正的 D-模。我们判断，这正是原文坚持 (11.1) 而不满足于「任意 F」的原因：\emph{条件 (11.1) 的功能不是让膜存在，而是让翻译有可积性}。\pB\anchor{src p.186}\pC\anchor{src p.198}

\paragraph{第五步：翻译与 S-对偶、't Hooft/Hecke 作用是否相容？}第 11 节不是孤立的一步：它必须与 S6、S7 拼起来才能得到几何 Langlands 的标准表述。

\emph{与 S-对偶相容。}S 把几何 Langlands 对偶等同于交换 \(\Psi=0\) 与 \(\Psi=\infty\) 的变换，在参数上作用为 \((G,\Psi)\mapsto(\Ldual G,-\Psi/n_g)\)（\(n_g=1\)，或对非单边群为 \(2,3\)）。\pA\anchor{src p.202} c.c. 膜由四维规范理论定义、不是 sigma 模型的临时构造，\pA\anchor{src p.187} 而 S-对偶把 \((A,B,A)\)-膜映到同型（第 12 节）；\pA\anchor{src p.208} 因此对 \(G\) 的翻译机制逐字重复即得 \(\Ldual G\) 一侧。\pA\anchor{src p.207} 更精确地，\(\Psi=h\)（\(h\) 为对偶 Coxeter 数）时 c.c. 弦代数退化为普通微分算子层；S-对偶与 \(T^h\) 合起来把 \(\MH(\Ldual G,C)\) 上的 zerobrane 映到 \(\MC\) 上的普通 D-模，这才是标准表述（细节见 \S\ref{subsec:twist}）。\pA\anchor{src p.207}

\emph{与 Hecke 相容。}S7（第 9--10 节）证明 't Hooft 线算子对应几何 Hecke 算子，S6（第 8 节）把 Wilson 算子对应到 electric eigenbrane；\pA\anchor{src p.185} 翻译命题把「膜」换成「D-模」，于是 't Hooft 作用变成 D-模上的 Hecke 作用，magnetic eigenbrane 的 Hecke 特征性变成「Hecke 特征层」的特征性。这条相容性是我们重建的（\pB）：'t Hooft 算子与 c.c. 膜都在同一个四维扭化 TFT 中定义，翻译所需数据（边界条件、拓扑超荷、开弦代数）都是该 TFT 的内在数据，故线算子作用与翻译应当交换；严格论证需对每个 Hecke 算子核验交换图，原文未逐条写出——这是全章最薄弱处之一（见 \S\ref{subsec:close}）。\pB\anchor{src p.185}\pC\anchor{src p.197}

\paragraph{例子 1：\(\varphi=0\) 的膜给出「定义环层的那个模」。}第一个例子是原文用来校准翻译的：取 \(B'\) 为 \(\varphi=0\) 处的膜、Chan-Paton 丛平凡，其支撑是 \(M=\MC\) 在 \(\MH\) 中的拷贝，在 \(K\) 与 \(J\) 中 Lagrangian、在 \(I\) 中复，所以是 \((B,A,A)\)-膜；第 10.4 节分析 't Hooft 算子 OPE 时用过它。\pA\anchor{src p.197} 计算 \((B_{c.c.},B')\) 弦的谱（原文约化到超对称量子力学，只留零模）：开弦态是 \(\mathcal N^{-1}\) 与「费米零模空间量子化所得丛」张成的丛的截面，而 \(\mathcal N\) 的联络限制到 \(\varphi=0\) 上为零，故可略去 \(\mathcal N^{-1}\)；费米零模空间就是 \(M\subset\MH\) 的切空间（法向费米子在弦两端有相反边界条件、无零模），量子化得 \(M\) 上旋量；\(M\) 复，spin 丛是 \(K_M^{1/2}\otimes(\bigoplus_j\Omega^{0,j}(M))\)，BRST 算子 \(Q\) 就是 \(K_M^{1/2}\) 值 \((0,j)\)-形式上的 \(\bar\partial\)。\pA\anchor{src p.198} 因此物理 \((B_{c.c.},B')\) 弦的层就是 \(K_M^{1/2}\) 的全纯截面层——正好是 \(\Dmod_{K_M^{1/2}}\) 的\emph{定义模}，这也说明环层为什么必须是 \(\Dmod_{K_M^{1/2}}\) 而非某个 \(\Dmod_L\)。\pA\anchor{src p.198}

原文还给了一个不依赖时间反演的论证：四维 SYM（\(M=\mathbb R\times I\times C\)，无涉时间的扭化）中膜是物理合理、酉的；沿时间方向除拓扑超荷 \(Q\) 外还有伴随 \(Q^\dagger\)，\(\{Q,Q^\dagger\}=H\)；charge conjugation 交换 \(Q\) 与 \(Q^\dagger\)，于是 Chan-Paton 丛必是 \(M\) 的 spin 丛、不乘额外线丛（[145] 在孤子量子化中首用此论证）。\pA\anchor{src p.198} 脚注 54 补充：\(j>0\) 的 \(\Omega^j\) 在小开集上对 BRST 上同调无贡献，可略去。

\paragraph{例子 2（几何 Langlands 的关键）：type F 膜 \(\mapsto\) D-模，\(G=U(1)\) 的完整计算。}第二个例子才是 GL 真正要用的：\(B'=\mathcal B_F\) 是卷在 Hitchin 纤维 \(F\) 上、带平坦 Chan-Paton 线丛的膜，即\term{type F brane}{type F 膜}；它是 magnetic eigenbrane，作为 \((A,B,A)\)-膜在 \(K\) 中是 A-膜，因此按命题 \ref{prop:translation} 给出 \(\Dmod_{K_M^{1/2}}\) 的模层。\pA\anchor{src p.199} 原文说，\(\S11.3\) 末会再给一个初等论证，把它转成 \(M\) 上普通微分算子 \(\Dmod\) 的模层；这一步其实就是 \(\Psi=h\) 时代数退化为普通 D（\S\ref{subsec:twist}），也是几何 Langlands 标准表述的最后一步。\pA\anchor{src p.199}

显式描述一般的 type F 膜很难，但 \(G=U(1)\) 可以完全算出来。设 \(G=U(1)\)：此时 \(A,\varphi\) 是实 1-形式，Hitchin 方程解耦；在复结构 \(I\) 中 \(\MH=\mathrm{Jac}(C)\times B\)，其中 \(B=H^0(C,\Omega^1)\)、\(\dim\mathrm{Jac}(C)=g_C\)，\(\MH\) 可等同于 \(\mathrm{Jac}(C)\) 的余切丛，Hitchin 纤维化就是到 \(B\) 的投影。\pA\anchor{src p.199} 取 \(p\in B\)，纤维为 \(F_p\simeq\mathrm{Jac}(C)\)，其上的 Chan-Paton 线丛 \(L\to F_p\) 拓扑平凡但带任意酉平坦联络。\pA\anchor{src p.200}

\paragraph{推导 C（要点）：\((B_{c.c.},B_{F_p})\) 弦给出平坦联络。}我们把原文 (11.17)--(11.25) 的要点写出来，并补出「为什么结果是平坦联络」。取 \(R\times I\) 上的世界面，左边界是 c.c. 膜、右边界是 \(\mathcal B_{F_p}\)；纤维 \(p\) 由 \(\varphi=\upsilon\) 定义，\(\upsilon=\upsilon_z\dd z+\upsilon_{\bar z}\dd\bar z\) 是实调和 1-形式。\pA\anchor{src p.200} 三项边界数据是 \(L\) 的平坦联络 \(\beta\)、\(\mathcal N\) 的联络 \(\alpha\)、辛势 \(\zeta\)（\(\omega=\dd\zeta\)）：右边界贡献 \(\int\Phi^*(\zeta-\ii\beta)\)、左边界贡献 \(-\int\Phi^*(\zeta-\ii\alpha)\)，相加得总作用量 (11.22)；它只依赖组合 \(\upsilon_z-a_z/\Im\tau\) 与 \(\upsilon_{\bar z}+a_{\bar z}/\Im\tau\)（11.23），不单独依赖 \(\upsilon\) 或 \(a\)。\pA\anchor{src p.200}\pA\anchor{src p.201} 第一项给出 \(\varphi_z\) 的量子化，第二项要求物理态被相应算子湮灭：
\begin{align}
\varphi_z&\ \longrightarrow\ -\frac{\ii\pi}{\Im\tau}\frac{\delta}{\delta A_z}+\frac{\ii\,\Im\tau}{2\pi}\upsilon_z-\frac{\ii}{2\pi}a_z
\quad\text{（}\mathrm{Jac}(C)\text{ 上平凡线丛的协变 }\bar\partial\text{）},\notag\\
0&=\Bigl(\frac{\delta}{\delta A_{\bar z}}-\frac{\ii\,\Im\tau}{2\pi}\upsilon_{\bar z}-\frac{\ii}{2\pi}a_{\bar z}\Bigr)\Psi
\quad\text{（协变 }\partial\text{）}.
\label{eq:phi-op}
\end{align}
\pA\anchor{src p.201} 现在补上原文没展开的一步（\pB）：这两者是一对 \((0,1)\)、\((1,0)\) 型协变导数，曲率由 \(\upsilon\) 的调和性（\(\bar\partial\upsilon=0\)）保证为零，故定义 \(\mathrm{Jac}(C)\) 上一个平坦联络，且当且仅当 \(\upsilon=0\) 时是酉的。因此
\begin{keyidea}
\((B_{c.c.},B_{F_p})\) 弦的量子化给出「\(\mathrm{Jac}(C)\) 上平凡线丛的复平坦联络」所对应的 D-模。
\end{keyidea}
\pA\anchor{src p.201} 意义：翻译得到的不是任意 D-模，而是由\emph{平坦联络}确定的 D-模，与 \S8 中 electric eigenbrane 由平坦局部系统确定的结构完全对称。原文随后指出 S-对偶把这族 A-膜与 \(\MH(U(1),C)\) 上的全体 zerobrane 对应，后者在复结构 \(J\) 中就是 \(C\) 上平坦 \(\CC^*\)-丛的模空间；于是阿贝尔情形下 S-对偶直接给出两侧平坦 \(\CC^*\) 联络的对应，后者可由 \(H_1(C,\ZZ)\cong H_1(\mathrm{Jac}(C),\ZZ)\)（Poincaré 对偶）直接看出。\pA\anchor{src p.201}

\begin{figure}[htbp]
\centering
\begin{tikzpicture}[font=\small,
  nd/.style={draw=kwblue!70, rounded corners=2pt, fill=kwboxbg, align=center,
             inner sep=5pt, text width=4.3cm, minimum height=1.2cm},
  arr/.style={-{Stealth[length=2.2mm]}, thick, draw=kwblue!70}]
  \node[nd] (brane) at (0,2.6) {\(K\) 中 A-膜\\（target \(\MH\)）};
  \node[nd] (dmod) at (-3.4,-1.6) {\(\Dmod_{K_M^{1/2}}\)-模\\（\(M=\MC\) 上，本章）};
  \node[nd] (nc) at (3.4,-1.6) {非交换形变 \(T^*\MC\) 上的 B-膜\\（半经典影子：\(T^*\MC\) 的相干层，[41]）};
  \draw[arr] (brane) -- node[left=2pt, align=right, font=\footnotesize] {(1) \((B_{c.c.},-)\)\\ 模结构} (dmod);
  \draw[arr, <->] (dmod) -- node[below=2pt, align=center, font=\footnotesize] {(2) symbol / 相伴分次（[41] 的等价）} (nc);
  \draw[arr] (brane) -- node[right=2pt, align=left, font=\footnotesize] {(3) 「A-膜 \(\leftrightarrow\) 非交换\\ B-膜」的特例} (nc);
\end{tikzpicture}
\caption{膜、D-模与非交换 B-膜三者的关系。箭头 (1) 是本章证明的翻译（单向：模结构）；箭头 (2) 是文献 [41] 的等价；箭头 (3) 是原文用来定位整件事的框架。原文并没有在本节证明 (2) 与 (1) 合起来构成范畴等价。}
\label{fig:triangle}
\end{figure}

\paragraph{排除错误直觉 4：\((A,B,A)\) 的「A」不是「这个膜不是 A-膜」。}读到「magnetic eigenbrane 是 \((A,B,A)\)-膜」，一个常见误解是：既然它标着 A 和 B 混合，它就不是 A-模型里的膜。错。\((A,B,A)\) 只是相对于三个复结构 \(I,J,K\) 的类型记录；本章（以及 \S11 全文）用的 A-模型是复结构 \(K\) 的那个，而 \(\mathcal B_F\) 在 \(K\) 中是 A-膜，所以它是这个 A-模型里堂堂正正的 A-膜。\pA\anchor{src p.208} 真正值得注意的是：同一个膜可以在不同复结构下有不同的类型；\((A,B,A)\) 说明它同时是 \(I\) 中 A-膜、\(J\) 中 B-膜。这种「跨复结构」的性质正是它能当 S-对偶载体的原因。\pA\anchor{src p.185}

\paragraph{本节收束。}把 \S11.2 的机制压成一句话：c.c. 膜的空间充满性给出「每个膜都是模」，(11.1) 给出「\(\hat N=I\) 可积」，时间反演给出「环层是 \(\Dmod_{K_M^{1/2}}\)」，三者合起来就是翻译。剩下两件事由 \S\ref{subsec:twist} 与 \S\ref{subsec:section12} 完成：把 \(\Psi=0\) 推广到任意 \(\Psi\)，以及用四维边界条件为这些膜提供独立来源。

\subsection{扭化微分算子：任意 \texorpdfstring{\(\Psi\)}{Psi} 下的翻译（\S11.3）}
\label{subsec:twist}

\paragraph{动机：为什么不能只停在 \(\Psi=0\)。}原文在 \S11.3 开头把几何 Langlands 对偶等同于「交换 \(\Psi=0\) 与 \(\Psi=\infty\) 的 S 变换」，并给出 S 在参数上的作用 \((G,\Psi)\mapsto(\Ldual G,-\Psi/n_g)\)；对有理 \(\Psi\) 最有趣（因为 \S6.2 的线算子在那里有丰富结构），而全 S-对偶群（由 \(T:\Psi\mapsto\Psi+1\) 与 \(S:\Psi\mapsto-\Psi/n_g\) 生成）把所有有理值互相联系起来。\pA\anchor{src p.202} 若翻译机制只在 \(\Psi=0\) 成立，就无法陈述 S-对偶本身。所以必须把 c.c. 膜推广到任意 \(\Psi\)。

为简化，原文取实 \(\Psi\)：此时 \(t=1\)、\(\theta\) 任意，且 \(\Psi=\theta/2\pi=\Re\tau\)。\pA\anchor{src p.202} 电与磁的理论这时更对称：两者都有 \(t=1\)，只是 \(\theta\) 互为倒数。把理论紧化在 \(C\) 上，得到 \(K\) 中的 A-模型，target 为 \(\MH(G,C)\)，但带 B-field \(B=-\omega_I\Re\tau\)（(4.19)）；S 变换把 \(\Re\tau\mapsto-\Re\tau/n_g\)、\(G\mapsto\Ldual G\)。\pA\anchor{src p.202}

立刻出现一个技术障碍：这个 B-field 沿 Hitchin 纤维不是平坦的（它是 \(\omega_I\) 的倍数，而 \(\omega_I\) 限制到纤维上非退化），于是对一般 \(\Psi\)，Hitchin 纤维不再是合法的 A-膜。\pA\anchor{src p.202} 这就解释了为什么第 11.3 节不能简单照抄第 11.1--11.2 节：A-膜的「候选对象」本身变了。

\paragraph{膜上的 B-field 与合法膜的条件。}膜卷在 \(Y\) 上时，规范变换 \(B\mapsto B-\dd\lambda\)、\(A\mapsto A+\lambda|_Y\) 使规范不变量是组合 \(F+B|_Y\)。\pA\anchor{src p.203} 当 \(X\) 辛、\(Y\) Lagrangian 时，\(B=0\) 的 A-膜要求 \(F=0\)；带 B-field 的一般条件是
\begin{equation}
F+B|_Y=0 .
\label{eq:F-plus-B}
\end{equation}
\pA\anchor{src p.203} 若 \([B|_Y/2\pi]\) 不是整数上同调类，则不存在满足 (11.29) 的 \(U(1)\) 规范场，也就没有单个 A-膜卷在 \(Y\) 上。于是：\(\theta\neq2\pi n\) 时 Hitchin 纤维不是合法 A-膜，不管 Chan-Paton 丛怎么取。\pA\anchor{src p.203} 一致地看，非整数 \(\Psi\) 的「naive T-对偶」——\(\MH(\Ldual G,C)\) 上的点——在 \(t=1\)、\(\theta\neq0\) 时也不是合法拓扑膜，因为点从来不是 A-膜（A-膜的支撑维数至少是中维）。\pA\anchor{src p.203}

原文的处理方式是：非整数 \(\Psi\) 下研究对象不是纤维上的 rank-one 膜；对有理 \(\Psi\)，纤维上有有限秩的 A-膜，可由 \(\Psi=0\) 沿对偶追踪得到。\pA\anchor{src p.203} 但无论怎样，S-对偶仍然把 \(\MH(G,C)\) 的 A-膜范畴与 \(\MH(\Ldual G,C)\) 的 A-膜范畴对应；而且对每个 \(\Psi\) 都存在 c.c. 膜的类比 \(\mathcal B_{c.c.}^{\Psi}\)，其开弦代数层是
\begin{equation}
\Dmod_{K_M^{1/2}\otimes L^{\Psi}}
\quad\text{（\(K_M^{1/2}\otimes L^{\Psi}\) 上的全纯微分算子层）},
\label{eq:twisted-alg}
\end{equation}
其中 \(L\to\MC\) 是行列式线丛（见式 (4.10)）；任意 A-膜 \(B'\) 的 \((B_{c.c.}^{\Psi},B')\) 弦层是这个代数的模。\pA\anchor{src p.203} 这就是扭化微分算子（\term{twisted differential operators}{扭化微分算子}）的来历：微分算子不再作用在函数上，而作用在线丛的复幂上。

\caveat{原文提醒：这些 twisted D-modules 与 \(\Psi=0\) 时一样，可以从四维规范理论构造，因此在代数几何上可以延拓到全 \(G\)-丛 stack \(\mathrm{Bun}_G(C)\) 上；S-对偶把 \(G\) 的 twisted D-模映到 \(\Ldual G\) 的、并交换 \(\Psi\mapsto-\Psi/n_g\)。\pA\anchor{src p.203}}

\paragraph{构造 \(\mathcal B_{c.c.}^{\Psi}\)。}空间充满 A-膜的规范场必须解
\begin{equation}
\bigl(\omega^{-1}(F+B)\bigr)^2=-1,\qquad \omega=(\Im\tau)\omega_K .
\label{eq:cc-psi}
\end{equation}
\pA\anchor{src p.204} 显然解为
\begin{equation}
F=\Im\tau\cos q\ \omega_J,\qquad \sin q=-\frac{\Re\tau}{\Im\tau},
\label{eq:F-psi}
\end{equation}
它定义 \(\mathcal B_{c.c.}^{\Psi}\)。\pA\anchor{src p.204} \(\Im\tau\) 可以任意大，此时 \(q\) 很小，有趣的效应线性于 \(\Psi=\Re\tau\sim q\)。

先看代数之外的结构：开弦层局部同构于 \(\MH\) 在复结构 \(I(\Psi)=\omega^{-1}(F+B)=I\cos q-J\sin q\) 中的全纯函数层，对应的全纯坐标是 \(A'_z=A_z-\ii\tan(q/2)\varphi_z\)、\(\varphi'_z=\varphi_z+\ii\tan(q/2)A_z\)。\pA\anchor{src p.204} Hitchin 第二纤维化 \(\pi_0:(A_z,\varphi_z)\mapsto A_z\)（(11.34)）除 \(\Psi=0\) 外对 \(I(\Psi)\) 并不全纯；把它改造成新联络 \(D'_z=D_z-\ii\tan(q/2)\varphi_z=\bar\partial_z+A_z-\ii\tan(q/2)\varphi_z\)，用 \(D'_z\) 定义全纯丛，得到 \(\pi_\Psi\)：对 \(I(\Psi)\) 全纯，\(\Psi=0\) 时退化为 \(\pi_0\)。\pA\anchor{src p.204} 脚注 57 指出 \(I(\Psi)=I_w\)，其中 \(w=\ii\tan(q/2)\)，正是 \S4.2 中那族等价的复结构。

纤维方向的结构用齐次对象编码：\(\varphi_z+\ii\tan(q/2)A_z\) 在规范变换下不齐次，齐次的是 \(D''_z=\varphi_z+\ii\tan(q/2)D_z\)；除以 \(\ii\tan(q/2)\) 后它是 \(E\) 上的 \(\bar\partial\)-联络（不除是为让 \(\Psi=0\) 也有意义）。原文说 \(D''_z\) 是 \(E'=\pi_\Psi(E,\varphi)\) 上的\emph{全纯 \(\lambda\)-联络}，\(\lambda=\ii\tan(q/2)\)，即线性映射 \(\nabla:\Gamma(E)\to\Gamma(E\otimes T^*X)\)，与 \(E\) 的 \(\bar\partial\) 交换且 \(\nabla(fs)=f\nabla s+\lambda\,\partial f\wedge s\)；\(\lambda=1\) 时是普通 \(\bar\partial\) 算子，\(\lambda=0\) 时是 \(\mathrm{End}(E)\otimes T^*X\) 的截面。由 Hitchin 方程，\(D''_z\) 与 \(D'_{\bar z}\) 交换，故 \(D''_z\) 确实定义 \(\lambda\)-联络；固定 \(E'\) 的 \(\lambda\)-联络空间是以 \(H^0(C,\mathrm{End}(E')\otimes K_C)\) 为模型的仿射空间，而后者正是 \(E'\) 处 \(M=\MC\) 的余切空间。\pA\anchor{src p.205}

于是得到原文的总结：忽略代数结构，\((B_{c.c.}^{\Psi},B_{c.c.}^{\Psi})\) 弦的层就是 \(M=\MC\) 的「扭化余切丛」上的全纯函数层；在复结构 \(I(\Psi)\) 中 \(\MH\) 看起来像 \(T^*\MC\) 的一个\emph{仿射形变}（纤维是仿射空间、底向量丛是 \(T^*\MC\)）。\pA\anchor{src p.205}

\paragraph{代数结构：为什么指数就是 \(\Psi\)。}代数结构与 \(\Psi=0\) 时同理：限制沿纤维多项式增长；局部量子化扭化余切丛与量子化普通余切丛无异，整体上只表现为修改 (11.16) 的转移函数。\pA\anchor{src p.206} 所以任意 \(\Psi\) 的 c.c. 弦层是 \(\MC\) 上「某个全纯线丛的幂」上的全纯微分算子层。\pA\anchor{src p.206}

\(\Psi\) 为整数时可以说得更精确。平移 \(\Psi\mapsto\Psi+n\) 把 \(B/2\pi\) 移动 \(L^{-n}\) 的第一 Chern 类（由 (4.19)：\(\Psi\mapsto\Psi+n\) 即 \(\theta\mapsto\theta+2\pi n\)，使 \(B\) 移动 \(-n\omega_I\)，而 \(\omega_I/2\pi=c_1(L)\)）。\pA\anchor{src p.206} 因此在 \(\Psi=n\) 处可以取膜 \(B_n\)：支撑在 \(\varphi=0\)、Chan-Paton 丛为 \(L^n\) 及其自然联络；它满足 (11.29)，故是 \(K\) 中的 A-膜。\pA\anchor{src p.206} 量子化 \((B_{c.c.}^n,B_n)\) 弦，物理态是
\begin{equation}
K_M^{1/2}\otimes L^n\otimes\Bigl(\bigoplus_{j=0}^{n}\Omega^{0,j}(M)\Bigr)
\label{eq:states-n}
\end{equation}
的截面；它是 \(\Dmod_{K_M^{1/2}\otimes L^n}\) 的模，而不是 \(\Dmod_{K_M^{1/2}\otimes L^m}\)（\(m\neq n\)）的模，所以代数层恰是 \(\Dmod_{K_M^{1/2}\otimes L^n}\)。\pA\anchor{src p.206}

\(\Psi\) 非整数时，不能指望有「态为 \(K_M^{1/2}\otimes L^{\Psi}\) 的截面」的膜（非整数时这样的模不存在），但已知代数层是 \(K_M^{1/2}\otimes L^{f(\Psi)}\) 上的微分算子层，只差确定 \(f\)。\pA\anchor{src p.206} 原文分两步：若 \(f\) 线性，则由整数处 \(f(n)=n\) 推出 \(f(\Psi)=\Psi\)；第二步用欧氏作用量（略去 \(\{Q,V\}\) 项）
\begin{equation}
S=\int_\Sigma\Phi^*(\omega-\ii F-\ii B)
=-\ii\,\Im\tau\int_\Sigma\Phi^*(\cos q\,\omega_J+\sin q\,\omega_I+\ii\,\omega_K)
\label{eq:S-euclid}
\end{equation}
证明线性。\(\Psi=0\) 时 \(q=0\)，被积 2-形式在 \(T^*\MC\) 上恰当，势（差数值因子）就是 \(p_\alpha\dd q^\alpha\)；\(\Psi\neq0\) 时写成「\(T^*\MC\) 上恰当形式 + 从 \(\MC\) 拉回的闭 2-形式」：
\begin{equation}
\Omega'=\frac{\Im\tau}{\pi}\,\delta\int_C\Tr(\varphi_z\,\delta A'_z)
-\ii\,\Re\tau\int_C\Tr(\delta A'_z\,\delta A'_{\bar z}),
\label{eq:Omega-prime}
\end{equation}
其中 \(A'_{\bar z}\) 是 \(A'_z\) 的负 Hermite 共轭。\pA\anchor{src p.207} 第二项是 \(\Re\tau\) 乘 \(\MC\) 上闭而非恰当的 \((1,1)\)-形式，是自然线丛 \(L\to M\) 曲率的倍数。\pA\anchor{src p.207} 计算 (11.16) 的转移函数要用 \(\Omega'\) 的逆作传播子；\([\hat p_\alpha,\hat p_\beta]\propto\hbar^2=1/(\Im\tau)^2\)，故只需算到 \(1/(\Im\tau)^2\) 阶，而在这一阶微扰论线性于 \(\Psi=\Re\tau\)，因为 \(\Omega'\) 的逆按 \(\sum_{k\ge0}a_k(\Re\tau)^k/(\Im\tau)^{k+1}\) 展开。\pA\anchor{src p.207} 证毕。

最后一步是常数：\(\MC\) 上的典范线丛同构于 \(L^{-2h}\)，其中 \(h\) 是 \(G\) 的对偶 Coxeter 数（[141]）。\pA\anchor{src p.207} 因此任意 \(\Psi\) 处 c.c. 弦的代数层是 \(L^{\Psi-h}\) 上的全纯微分算子层。特别地，令
\begin{equation}
\Psi=h,
\label{eq:Psi-h}
\end{equation}
则代数层就是\emph{普通}微分算子层（作用在函数上）。\pA\anchor{src p.207}

\begin{keyidea}
原文的收束句可以逐字翻译为：把 S-对偶与平移 \(\Psi\mapsto\Psi+h\) 的变换 \(T^h\) 合起来，就能把 \(\MH(\Ldual G,C)\) 上的 zerobrane 映到 \(\MC\) 上的普通（未扭化）D-模。\pA\anchor{src p.207} 这就是几何 Langlands 的标准表述；对某些 \(G\)，考虑 \(C\) 的 spin 结构依赖后它还是更精确的表述（[5,26]）。\pA\anchor{src p.207}
\end{keyidea}

\begin{figure}[htbp]
\centering
\begin{tikzpicture}[font=\small,
  arr/.style={-{Stealth[length=2.2mm]}, thick},
  pt/.style={circle, fill=kwblue!70, inner sep=1.6pt},
  lab/.style={align=center, font=\footnotesize}]
  \draw[kwgray!60, thick, -{Stealth[length=2.2mm]}] (-0.6,0) -- (11.4,0);
  \node[lab, below] at (11.2,-0.25) {\(\Psi\)};
  \node[pt] (p0) at (0,0) {};
  \node[pt] (pn) at (3.2,0) {};
  \node[pt] (ph) at (6.0,0) {};
  \node[pt] (pg) at (9.4,0) {};
  \node[lab, above=3pt] at (p0) {\(\Psi=0\)};
  \node[lab, below=4pt, text width=4.6cm] at (p0) {代数 \(=\Dmod_{K_M^{1/2}}\)\\（\S11.2 的结果）};
  \node[lab, above=3pt] at (pn) {\(\Psi=n\in\ZZ\)};
  \node[lab, below=4pt, text width=4.6cm] at (pn) {代数 \(=\Dmod_{K_M^{1/2}\otimes L^n}\)\\（\(B_n\) 带 Chan-Paton \(L^n\)）};
  \node[lab, above=3pt] at (ph) {\(\Psi=h\)};
  \node[lab, below=4pt, text width=4.2cm] at (ph) {代数 \(=\Dmod\)（普通）\\标准几何 Langlands};
  \node[lab, above=3pt] at (pg) {一般 \(\Psi\)};
  \node[lab, below=4pt, text width=4.0cm] at (pg) {代数 \(=\Dmod_{K_M^{1/2}\otimes L^{\Psi}}\)};
  \draw[arr, thick, draw=kwred!70] (p0) to[bend left=18] node[above=2pt, font=\scriptsize] {\(T\)：\(\Psi\mapsto\Psi+1\)} (pn);
  \draw[arr, thick, draw=kwred!70] (pn) to[bend left=18] node[above=2pt, font=\scriptsize] {\(T\)} (pg);
  \draw[arr, thick, draw=kwgreen!70] (pg) to[bend left=22] node[below=2pt, font=\scriptsize] {\(T^{h-\Psi}\) 把一般值移到 \(\Psi=h\)} (ph);
\end{tikzpicture}
\caption{\(\Psi\) 轴上的 c.c. 代数。\(\Psi=0\) 是本章第 11.2 节算出的扭化微分算子层；整数 \(\Psi=n\) 处可用带 Chan-Paton 线丛 \(L^n\) 的膜 \(B_n\) 来探测，指数与 \(\Psi\) 相同；\(\Psi=h\) 处退化为普通微分算子层，这一步配合 S-对偶给出标准几何 Langlands 表述。}
\label{fig:ladder}
\end{figure}

\paragraph{非交换结构：Kontsevich 与 Seiberg--Witten。}原文在 \S11 开头就把 D-模定位为「\(\MC\) 余切丛的非交换形变上的 B-膜」，并引用 [41] 的等价作为一般框架。\pA\anchor{src p.185} 第 11.2 节的具体计算则给出了这个非交换代数的一个显式模型：c.c. 弦代数就是 \(\Dmod_{K_M^{1/2}}\)，而 \(\Dmod_{K_M^{1/2}}\) 正是「\(\MH\) 在复结构 \(I\) 下的全纯函数代数」的形变量子化——形变参数来自 \((g+F)^{-1}\) 的反对称部分（Yang--Mills 文献 [137] 的 \(\theta\)），在本问题中正比于 \(\omega_J^{-1}\)，而 \(\omega_J=\Re\Omega_I\)。\pA\anchor{src p.189} 把 \(\omega_J^{-1}\) 换成 \(\Omega_I^{-1}\) 后，全纯函数的乘法就按 \(\Omega_I\) 的 Poisson 括号形变。\pA\anchor{src p.189}

与 \term{Kontsevich quantization}{Kontsevich 量子化} 的关系：复辛流形（这里是 \(\MH\) 的开集，复结构 \(I\)、全纯辛形式 \(\Omega_I\)）上的形变量子化总存在，Kontsevich 定理保证这一点，而 \(\Dmod_{K_M^{1/2}}\) 是它的具体实现——把纤维上的多项式函数量子化为微分算子就是标准符号演算。我们判断，原文的「非交换形变」指的是\emph{\(\MH\) 的函数代数}（半经典地看是 \(T^*\MC\) 的函数代数）的形变，而不是 \(\MC\) 本身；因此把 D-模说成「\(\MC\) 上非交换空间的层」并不准确：非交换只出现在余切方向。\pB\anchor{src p.185}\pC\anchor{src p.189}

与 \term{Seiberg--Witten}{Seiberg--Witten} 非交换几何的关系更微妙：那里的非交换参数同样是 \((g+F)^{-1}\) 的反对称部分，但通常要求 \(|F|\gg|g|\)，本题恰好 \(|F|=|g|\)。原文的处理是：限制到拓扑可观测量后弦激发被消去，形变代数仍可定义，但只是「拓扑影子」。这也解释了为什么第 11 节要算到二阶微扰并在整体上补转移函数：它要的不是关于非交换几何的一般定理，而是这个具体代数层的精确模型——这正是「典范」的另一个含义：形变由 \(\Omega_I\) 唯一决定。\pB\anchor{src p.189}\pC\anchor{src p.191}

\caveat{脚注 58 是本节最容易被忽略的技术点：在 \(\Psi\neq0\) 时，对固定复结构 \(I(\Psi)\) 存在两种相容的代数结构——\(G\)-丛带 \(\lambda\)-联络的模空间，以及稳定同态 \(\pi_1(C)\to G_{\CC}\) 的模空间；原文在此必须选择前者（沿纤维多项式增长的函数），并说明这是全文唯一一次必须在两种代数结构间做选择。\pA\anchor{src p.206}}

\paragraph{排除错误直觉 5：「\(\Psi\neq0\) 只是把 D-模扭一下，几何不变」。}错。\(\Psi\neq0\) 时至少有三件本质变化：第一，A-模型的 B-field 非零，底层的复结构从 \(I\) 变成 \(I(\Psi)=I\cos q-J\sin q\)；第二，c.c. 膜仍是空间充满的，但它所在的 A-模型、它的高阶结构都随之改变；第三，代数层从 \(\Dmod_{K_M^{1/2}}\) 变成 \(\Dmod_{K_M^{1/2}\otimes L^{\Psi}}\)，而非整数 \(\Psi\) 时连线丛 \(L^{\Psi}\) 都不存在，只有扭化微分算子层作为抽象层存在。\pA\anchor{src p.204}\pA\anchor{src p.206} 所以「扭化」不是装饰，它是 S-对偶群作用在几何上的真实体现。

\paragraph{排除错误直觉 6：「\(\Psi=h\) 的普通 D-模是 \(\Psi=0\) 时扭化 D-模的特例」。}普通 D-模并不是「同一个理论里把扭化参数调到 0」；它是另一个 \(\Psi\) 值处的代数，并且要配合对偶群的一侧来读：原文的陈述是「S-对偶 \(+\) \(T^h\) 把 \(\MH(\Ldual G,C)\) 上的 zerobrane 映到 \(\MC\) 上的普通 D-模」。\pA\anchor{src p.207} 换句话说，标准表述的「普通 D-模」是膜侧 zerobrane、\(\Psi=h\)、以及群为对偶群三者一起给出的；只盯着 \(\Psi=0\) 的代数会看不出这一点。

\subsection{第 12 节：检查清单——三类膜的规范理论边界条件}
\label{subsec:section12}

\paragraph{动机：为什么第 11 节之后还需要第 12 节。}第 11 节的翻译是在 target 空间 \(\MH\) 的 sigma 模型里做的，但原文在 \S11 开头就承认：几何 Langlands 的 D-模必须做在 \(\mathrm{Bun}_G(C)\) 这个 stack 上，只做在稳定丛模空间不够。\pA\anchor{src p.185} 第 12 节的目标正是补上这一点：说明重要的膜可以在底层四维规范理论中直接定义，而不只是在 target 为 \(\MH\) 的 sigma 模型里定义。\pA\anchor{src p.208} 因此第 12 节的作用是\emph{自洽性检查}：它用四维边界条件重新造出第 11 节用到的三类膜，验证 \((A,B,A)\)、\((B,A,A)\)、\((B,B,B)\) 这些类型不是 sigma 模型近似的产物，并把 c.c. 膜的定义从「\(\MH\) 的光滑部分」提升到「整个 \(G\)-丛 stack」。\pA\anchor{src p.187}

\caveat{为什么这件事不能省？因为第 11 节的公式 (11.2) 只在 \(\MH\) 光滑部分成立，而 c.c. 膜的开弦代数又决定了整章翻译。若 c.c. 膜只是 \(\MH\) 上的一个对象，那么「D-模做在 \(\mathrm{Bun}_G(C)\) 上」这一步就无从谈起。第 12.4 节的边界条件恰好给出不依赖 sigma 模型近似的定义。\pA\anchor{src p.187}}

\paragraph{第 12.1 节的一般方法：从边界项到超对称边界条件。}原文的分析套路是纯经典的：先算玻色子作用量的边界项，再算费米子的，最后要求超流法向分量为零。规范场与费米子作用量及其边界项为
\begin{align}
I_B&=\frac{1}{2e^2}\int\dd^n x\sqrt g\,\Tr F_{IJ}F^{IJ},\qquad
(\delta I_B)_{\mathrm{bdry}}=\frac{2}{e^2}\int\dd^{n-1}x\sqrt g\,\Tr\Bigl(g^{IJ}\sum_{I,J\neq1}\delta A_I F_{1J}\Bigr),\label{eq:IB}\\
I_F&=-\frac{\ii}{e^2}\int\dd^n x\sqrt g\,\Tr\bar\lambda\Gamma^I D_I\lambda,\qquad
(\delta I_F)_{\mathrm{bdry}}=-\frac{\ii}{e^2}\int\dd^{n-1}x\sqrt g\,\Tr\bar\lambda\Gamma^1\delta\lambda .\label{eq:IF}
\end{align}
\pA\anchor{src p.210} 令边界项为零有两条基本途径：令 \(\delta A_I=0\)（对 \(A_I\) 加 \term{Dirichlet boundary condition}{Dirichlet 边界条件}），或在 \(A_1=0\) 规范下令 \(F_{I1}=0\)（\term{Neumann boundary condition}{Neumann 边界条件}的规范不变版本）；更一般地可令二者的某个最大线性组合为零，称为 \term{hyperbolic boundary condition}{双曲边界条件}。费米子侧则令 \(\lambda\) 的一半分量为零，使 \(\bar\lambda\Gamma^1\delta\lambda=0\)。\pA\anchor{src p.210} 超流与超对称条件为
\begin{equation}
J_I=\frac12\Gamma^{KL}F_{KL}\Gamma_I\lambda,\qquad
\epsilon J_1=0
\ \Longleftrightarrow\
\delta\lambda=\frac12\Gamma^{IJ}F_{IJ}\epsilon=0
\quad\text{（在边界上）}.
\label{eq:supercurrent}
\end{equation}
\pA\anchor{src p.211} 实际操作是：先假定费米子的简单边界条件，推出玻色子边界条件，再验证 \(\epsilon J_1=0\)。原文强调：若边界条件自发破缺超对称就没有意义；而玻色子的 Dirichlet 条件一般与未破缺超对称不相容，所以有意义的构造是「固定 \(F_{IJ}\) 的边界值，并使之与 \(\delta\lambda=0\) 相容」。\pA\anchor{src p.211}

还有全局的计数：膜保两个拓扑超荷（加伴随共四个超对称），从四维 \(N=4\) 理论看是 \(1/4\) BPS，从相对超对称已减半的二维 sigma 模型看是 \(1/2\) BPS；未破缺旋量由扭化条件
\begin{equation}
\Gamma^{2367}\epsilon=\epsilon
\label{eq:chirality}
\end{equation}
刻画（同原文式 (5.3)）。\pA\anchor{src p.209}

\paragraph{\((B,B,B)\)：zerobrane 及其第二个构造。}取 \(T=\Gamma^{04}\)，费米子边界条件为 \(T\lambda|=\lambda|\)（\(|\) 表示限制到 \(x^1=0\)）。\pA\anchor{src p.212} 要求 \((T-1)\delta\lambda=0\)，代入 \(\delta\lambda=\frac12\Gamma^{IJ}F_{IJ}\epsilon\) 并用 \(T\epsilon=\epsilon\)，得 \([T,\Gamma^{IJ}F_{IJ}]\epsilon=0\)，即
\begin{equation}
\sum_{I\neq0,4}\bigl(\Gamma^{0I}F_{4I}-\Gamma^{4I}F_{0I}\bigr)\epsilon=0
\ \xRightarrow{\ \Gamma^0\Gamma^4\epsilon=\epsilon\ }\
F_{0I}+F_{4I}=0\quad(I\neq0,4).
\label{eq:BB-cond}
\end{equation}
\pA\anchor{src p.212} 取双曲边界条件：规范下设 \(A_0+A_4=0\)，对 \(I\neq0,1,4\) 的 \(A_I\) 加 Dirichlet，则 (12.12) 就是 \(A_0+A_4\) 的规范不变 Neumann；超流法向条件 \(\epsilon\Gamma^{MN}F_{MN}\Gamma^1\lambda=0\) 与玻色变分相容性都自动满足。\pA\anchor{src p.213}

还需要存在超对称态：若对所有 \(\epsilon\) 要求 \(\frac12\Gamma^{IJ}F_{IJ}\epsilon=0\) 则 \(F=0\)；但 \(\epsilon\) 受 \(\Gamma^{2367}\epsilon=\epsilon\) 限制，故只要求 \(F\) 在 \(x^2,x^3,x^6,x^7\) 张成的子空间上自对偶，即 Hitchin 方程
\begin{equation}
F_{23}-F_{67}=0,\qquad F_{27}-F_{36}=0,\qquad F_{26}-F_{73}=0 .
\label{eq:hitchin-12}
\end{equation}
\pA\anchor{src p.214} 边界值因此定义 \(\MH\) 中一个点 \(p\)。再用「除非 \(I,J\) 都在 \(R=\{2,3,6,7\}\) 中，\(F_{IJ}|=0\)」：对 \(J=5,8,9\)，\(A_J|\) 在 \(C\) 上协变常数且与 Higgs 场（\(\varphi_2=A_6,\varphi_3=A_7\)）交换；\(p\) 是不可约 Higgs 丛的光滑点时这迫使 \(A_J|=0\)，进而 \(F_{IJ}|=0\)（\(I,J\notin R\)），边界条件唯一（\(p\) 可约时只能确定 \(A_5,A_8,A_9\) 生成互交换对称）。结论：支撑在 \(p\) 的 zerobrane，最基本的 \((B,B,B)\)-膜。\pA\anchor{src p.214}

第二个构造：把 \(T\) 换成 \(\tilde T=T\Gamma^{2367}=\Gamma^{023467}\)，取 \(\tilde T\lambda|=-\lambda|\)。\pA\anchor{src p.214} 超对称条件加 \(\Gamma^{2367}\epsilon=\epsilon\) 后化为 (12.17)、(12.18) 两组：前者由 \(A_I\)（\(I\in\{2,3,6,7\}\)）的 Dirichlet、Hitchin 方程与 \(A_4|=0\) 实现；后者不能读成「\(F_{IJ}|=0\)」（会超定）——由十维手征性 \(\hat\Gamma\epsilon=\epsilon\) 与 \(\tilde T\epsilon=\epsilon\) 得 \(\Gamma^{1589}\epsilon=-\epsilon\)，于是它等价于 1589 平面中边界值的\emph{反对偶}条件
\begin{equation}
F_{15}|+F_{89}|=0,\qquad F_{18}|+F_{95}|=0,\qquad F_{19}|+F_{58}|=0,
\label{eq:BB2-3}
\end{equation}
这是修改版 Neumann 条件。\pA\anchor{src p.215} 其余部分同第一个构造（\(A_2,A_3,A_6,A_7\) 由 Hitchin 解加 Dirichlet），故 zerobrane 仍支撑在 \(p\)；差别只在 \(A_4,A_5,A_8,A_9\) 的边界条件以及平坦时空中超对称数少一半。\pA\anchor{src p.215}

\paragraph{两个构造等价吗？}当 \(p\) 光滑不可约时，\(A_4,A_5,A_8,A_9\) 没有零模，低能理论中无论微观边界条件如何它们都消失，故两个构造在光滑 \(p\) 处应当等价；\(p\) 可约（稳定子正秩）时这些场有零模、生成 Higgs 丛的对称，两个构造没有理由等价。\pA\anchor{src p.216} 这正是原文反复提醒的「提升不唯一」：同一张 2d 膜可以有不同的 4d 边界条件实现，它们只在奇点附近才显出差别。\pA\anchor{src p.208}

\paragraph{\((B,A,A)\)：S-对偶像与 Hitchin 纤维。}第二个类型由 \((B,B,B)\) 经 S-对偶得到。原文说：具体膜如何变换一般不清楚，但未破缺超对称可以变换；由式 (2.25)，在 \(\Re\tau=0\) 时 S 作用为 \(\epsilon_l\mapsto\epsilon_l\ee^{-\ii\pi/4}\)、\(\epsilon_r\mapsto\epsilon_r\ee^{\ii\pi/4}\)，等价于 \(\epsilon\mapsto\epsilon'=\frac{1+\hat\Gamma}{\sqrt2}\epsilon\)（\(\hat\Gamma=\Gamma^0\Gamma^1\Gamma^2\Gamma^3\)）。\pA\anchor{src p.216} 原来的 \(T\epsilon=\epsilon\) 变成 \(T'\epsilon'=\epsilon'\)，其中
\begin{equation}
T'=\frac{1+\hat\Gamma}{\sqrt2}\,T\,\frac{1-\hat\Gamma}{\sqrt2};
\qquad T=\Gamma^0\Gamma^4\ \Longrightarrow\ T'=\Gamma^1\Gamma^2\Gamma^3\Gamma^4 .
\label{eq:Tprime}
\end{equation}
\pA\anchor{src p.216} 取 \(T'\lambda|=-\lambda|\)，超对称条件 \(\{\Gamma^{IJ}F_{IJ},T'\}\epsilon=0\) 即
\begin{equation}
\Bigl(\sum_{I,J\in\{1,2,3,4\}}+\sum_{I,J\notin\{1,2,3,4\}}\Bigr)\Gamma^{IJ}F_{IJ}\,\epsilon=0 .
\label{eq:BAA-cond}
\end{equation}
\pA\anchor{src p.216} 由它得：对 \(I>4\)，\(F_{0I}|=0\)，即 \(A_I\) 加 Dirichlet；置 \(A_I|=0\)（\(I=5,8,9\)），寻找非零的「Higgs」分量 \(\varphi_2=A_6,\varphi_3=A_7\)。若条件对所有满足 \(\Gamma^{1234}\epsilon=\epsilon\)（不要求 \(\Gamma^{2367}\epsilon=\epsilon\)）的 \(\epsilon\) 成立，则 \([\varphi_2,\varphi_3]=F_{67}=0\)：给定 \((\varphi_2,\varphi_3)\) 的规范不变内容等价于给定 \(\varphi_z=\frac12(\varphi_2-\ii\varphi_3)\) 的特征多项式，于是边界条件就是「特征多项式固定」的 Dirichlet。\pA\anchor{src p.217} \(A_2,A_3,A_4\) 满足修改的 Neumann 条件 \(F_{12}|-F_{34}|=F_{13}|-F_{42}|=F_{14}|-F_{23}|=0\)；\(\Gamma^{2367}\epsilon=\epsilon\) 的不破缺又要求 \(F_{26}|-F_{73}|=F_{27}|-F_{36}|=0\)，即 \(D_{\bar z}\varphi_z=0\)，原文称「复」Hitchin 方程。于是决定全纯结构的 \(A_z\) 可涨落、\(\varphi_z\) 的特征多项式固定，故膜支撑在 Hitchin 映射 \(\pi:\MH\to B\) 的纤维上，正是原型 \((B,A,A)\)-膜。\pA\anchor{src p.217}

原文有一个诚实的注记：经典分析中出现的 \([\varphi_2,\varphi_3]=0\) 并非 Hitchin 方程之一（后者应为 \([\varphi_2,\varphi_3]-F_{23}=0\)），多出的项大概是重整化群流向红外时产生的量子修正。\pA\anchor{src p.217} 二维极限下，S-对偶把 Hitchin 纤维上的 \((B,A,A)\)-膜映到 \((B,B,B)\)-zerobrane；阿贝尔情形可直接验证。\pA\anchor{src p.217}

\((B,A,A)\) 还有第二个构造：\(T''=T'\Gamma^{2367}=-\Gamma^{1467}\)，取 \(T''\epsilon=\epsilon\)、\(T''\lambda|=\lambda|\)，则对 \(A_6=\varphi_2\)、\(A_7=\varphi_3\) 加 Dirichlet、对 \(A_2,A_3\) 加 Neumann。\pA\anchor{src p.218} Hitchin 纤维 \(F\) 上应当有许多膜（对应 \(F\) 上的平坦线丛），但时间反演对称使上述两个构造都带平凡 Chan-Paton（至多相差二阶）；原文希望修改规范构造使丛可变。\pA\anchor{src p.218}

\paragraph{\((A,B,A)\)：c.c. 膜的规范理论边界条件。}这是第 12 节最重要的构造，因为它要认出的正是第 11 节的 c.c. 膜。原文先引入缩写 \(\Gamma_{\pm I\pm J}=(\pm\Gamma_I\pm\Gamma_J)/\sqrt2\)、\(F_{\pm I\pm J,K}=(\pm F_{IK}\pm F_{JK})/\sqrt2\)，再定义（原文在 \S12.4 中重新使用了 \(T'\) 这个记号；为避免与 (12.23) 的 \(T'\) 混淆，本节记它为 \(\mathcal T\)）
\begin{equation}
\mathcal T=-\Gamma_{15}\,\Gamma_{2-7}\,\Gamma_{3+6}.
\label{eq:T-ABA}
\end{equation}
\pA\anchor{src p.218} 这个 \(\mathcal T\) 有一个关键性质：它在 \(x^2\)-\(x^3\) 平面与 \(x^6\)-\(x^7\) 平面\emph{同时}旋转下不变。因此在扭化理论中 \(\mathcal T\) 不依赖正交坐标 \(x^2,x^3\) 的选取；在 \(M=\RR^2_+\times C\) 上它只依赖 \(C\) 的度规而不依赖坐标，于是所造的膜只依赖度规的共形类，即只依赖 \(C\) 的复结构。\pA\anchor{src p.218}

分两种情形。第一种取 \(\mathcal T\lambda|=\lambda|\)。超对称条件 \([\Gamma^{IJ}F_{IJ},\mathcal T]\epsilon=0\) 给出：若 \(I,J\) 之一在 \(\mathrm{span}\{1,5,2-7,3+6\}\) 中而另一个与之正交，则 \(F_{IJ}=0\)。特别地，\(F_{0I}|=0\)（\(I=5,2-7,3+6\)）给出 Dirichlet，\(F_{1I}=0\)（\(I=2+7,3-6,4,8,9\)）给出 Neumann。\pA\anchor{src p.218} 所得的膜是：\(A_2-\varphi_3\)、\(A_3+\varphi_2\) 有指定边界值，\(A_2+\varphi_3\)、\(A_3-\varphi_2\) 自由涨落。\pA\anchor{src p.219}

用复结构 \(J\) 中的全纯坐标来读这个条件。直接算：
\begin{align}
2(A_z+\ii\varphi_z)&=(A_2-\ii A_3)+\ii(\varphi_2-\ii\varphi_3)
=(A_2+\varphi_3)-\ii(A_3-\varphi_2),
\label{eq:J-coord1}\\
2(A_{\bar z}+\ii\varphi_{\bar z})&=(A_2+\ii A_3)+\ii(\varphi_2+\ii\varphi_3)
=(A_2-\varphi_3)+\ii(A_3+\varphi_2).
\label{eq:J-coord2}
\end{align}
\pA\anchor{src p.219} 所以条件「\(A_2-\varphi_3\) 与 \(A_3+\varphi_2\) 固定」恰好是说复合场 \(\mathcal A_z=A_z+\ii\varphi_z\) 固定，而 \(\mathcal A_{\bar z}\) 自由。于是 \(M_H\) 的子流形 \(W'\) 在复结构 \(J\) 中全纯；而全纯辛形式 \(\Omega_J=-\frac{\ii}{4\pi}\int_C\Tr\delta A\wedge\delta A\) 在 \(W'\) 上为零（\(W'\) 中维且 \(\delta A_z=0\)），所以 \(W'\) 是 \(J\) 的复 Lagrangian 子流形。\pA\anchor{src p.219} 因此 \(W'\) 上平凡 Chan-Paton 的膜是 \((A,B,A)\)-膜；\(W'\) 的几何意义是「带指定全纯结构的平坦 \(G_{\CC}\)-丛」。\pA\anchor{src p.219}

第二种情形取 \(\mathcal T\lambda|=-\lambda|\)。此时 \([\Gamma^{IJ}F_{IJ},\mathcal T]\epsilon=0\) 给出：对与 \(\mathrm{span}\{1,5,2-7,3+6\}\) 垂直的 \(I,J\) 有 \(F_{IJ}=0\)；那里 \(F_{0I}=0\) 意味着 \(A_2+\varphi_3\)、\(A_3-\varphi_2\) 加 Dirichlet；而 \(F\) 在 \(\mathrm{span}\{1,5,2-7,3+6\}\) 上反对偶，给出一组修改的 Neumann 条件
\begin{align}
F_{15}+F_{2-7,3+6}&=0,\notag\\
F_{1,2-7}+F_{3+6,5}&=0,\label{eq:ABA-neu}\\
F_{1,3+6}+F_{5,2-7}&=0.&\notag
\end{align}
\pA\anchor{src p.219} 这第二种膜中 \(A_{2+7},A_{3-6}\) 固定而 \(A_{2-7},A_{3+6}\) 自由，恰好是 \(A_z\) 与 \(A_{\bar z}\) 的角色互换，所以它支撑在 \(W''\)——平坦 \(G_{\CC}\)-丛加上固定的\emph{反全纯}结构——上，仍是 \((A,B,A)\)。\pA\anchor{src p.219} 把 \(\varphi_2,\varphi_3\) 同时变号还可以得到另一族（\(\mathcal T''=-\Gamma_{15}\Gamma_{2+7}\Gamma_{3-6}\)）；两族只差在远离奇点处 massive 的模式，可能在奇点处不同。\pA\anchor{src p.219}

\paragraph{第 12.4 节的收束：这个 \((A,B,A)\) 膜就是 c.c. 膜。}怎么知道上面造出的空间充满 \((A,B,A)\) 膜就是第 11 节的 \(\mathcal B_{c.c.}\)？原文用「混合边界条件的曲率」来认它。膜上非平坦规范场的一般边界条件形如 \(f_{MN}\partial_0X^N+g_{MN}\partial_1X^N=0\)，其中 \(f\) 是 target 空间规范场（brane 上的场强）、\(g\) 是度规；\(f=0\) 时退化为 Neumann。\pA\anchor{src p.221} 把 S-对偶作用到上面的 \((A,B,A)\) 膜后，原文算出边界条件（式 (12.31)、(12.32)）：
\begin{align}
&F_{15}+F_{26}+F_{37}=0,\quad
F_{12}+F_{06}-F_{56}=0,\quad
F_{13}+F_{07}-F_{57}=0,\quad F_{0J}=0\ (J=4,8,9);\label{eq:ABA-S}\\
&D_0\varphi_2+\partial_1A_2=0,\quad
D_0\varphi_3+\partial_1A_3=0,\quad
F_{02}-\partial_1\varphi_2=0,\quad
F_{03}-\partial_1\varphi_3=0,\label{eq:ABA-32}
\end{align}
后者是置 massive 的 \(A_4,A_8,A_9\) 为零、取 \(A_1=0\) 规范之后得到的。\pA\anchor{src p.220} 与混合边界条件比较，可知这里的 target 规范场是 \(f=\frac{1}{2\pi}\int_C|\dd^2z|\,\Tr(\delta\varphi_z\wedge\delta A_z+\delta\varphi_{\bar z}\wedge\delta A_{\bar z})=\omega_J\)。\pA\anchor{src p.221} 而「target 规范场为 \(\omega_J\) 的空间充满 \((A,B,A)\) 膜」正是第 11.1 节 c.c. 膜的定义。至此完成闭环：第 11 节用 sigma 模型语言定义的 \(\mathcal B_{c.c.}\)，在第 12 节用四维边界条件重新造了出来。

\paragraph{c.c. 膜的边界可观测量：与第 11.2 节的代数对上。}原文最后把第 11.2 节的开弦代数与这里独立得到的边界条件对应起来：考虑插入在边界点 \(P\) 的局域算子 \(\mathcal O\)（边界形如 \(\RR\times C\)）；\(C\) 的度规模 \(\{Q,\dots\}\) 无关，故可设 \(C\) 在 \(p\) 附近平坦；按维数与「绕 \(p\) 旋转的自旋」分类，自然维数为 \(A,\varphi,\psi,\bar\psi,e\) 为 1、\(\chi,\eta,\bar\eta\) 为 2。\pA\anchor{src p.221} 维数 \(n\)、自旋 \(n\) 的 \(Q\) 不变算子只能由 \(\varphi_z,\psi_z,\bar\psi_z\) 与 \(D_z\) 构成，而 \(Q\) 不变要求 \(D_z\) 不出现；c.c. 膜边界条件 \(\hat T\psi|=-\psi|\) 给出 \(\psi_2|=-\ii\bar\psi_3|\)、\(\psi_3|=\ii\bar\psi_2|\)，故 \(\psi_z|=\bar\psi_z|\)。于是这类可观测量只是 \(\varphi_z,\psi_z\) 的函数、不含导数：\(K=0\) 时如 \(\mathcal O=\Tr\varphi_z^n\)，\(K=1\) 时如 \(\mathcal O'=\Tr\varphi_z^{n-1}\psi_z\)。\pA\anchor{src p.221}

量子修正不会让 \(\mathcal O\) 与 \(\mathcal O'\) 配对：时间反演 \(T\) 反 \(x^0\)、对其它坐标平凡，是 c.c. 膜边界条件的对称且与拓扑超荷交换（\(t=1\) 是 \(T\) 的不动点）；\(\varphi_z\) 在 \(T\) 下为奇而 \(\psi_z|\) 为偶，故不能配对。一般地，自旋 \(n\)、ghost number \(k\) 的经典 \(Q\) 不变算子由 \(F(\varphi_z,\psi_z)\) 决定、在 \(T\) 下变 \((-1)^{n-k}\)；\(Q\) 保持 \(n\) 而把 \(k\) 加 1，翻转 \(T\) 本征值，故不可能。所以 \(Q\) 的上同调等于经典结果。\pA\anchor{src p.222} 与第 11.2 节对照：那边说「开弦代数是什么」（\(\Dmod_{K_M^{1/2}}\)），这里说「开弦态在四维边界条件里如何实现」；结论一致，这正是第 12 节作为检查清单的价值。\pB\anchor{src p.221}

\paragraph{为什么只有这些类型：脚注 24 与三复结构表。}把三类膜的 A/B 类型按复结构 \(I,J,K\) 列出来，就得到原文的「类型表」。图~\ref{fig:types} 把它画成一张表，并标出 S-对偶的作用。

\begin{figure}[htbp]
\centering
\begin{tikzpicture}[font=\small,
  hdr/.style={font=\footnotesize\bfseries, text=kwblue, align=center},
  cell/.style={align=center, inner sep=2pt, minimum width=1.9cm, minimum height=0.62cm},
  rlab/.style={font=\footnotesize, text=kwgray, align=right, inner sep=3pt}]
  % header row
  \node[hdr] at (2.6,2.55) {zerobrane};
  \node[hdr] at (5.0,2.55) {Hitchin 纤维};
  \node[hdr] at (7.4,2.55) {c.c. 膜};
  \node[hdr, text=kwgray] at (2.6,3.1) {\((B,B,B)\)};
  \node[hdr, text=kwgray] at (5.0,3.1) {\((B,A,A)\)};
  \node[hdr, text=kwgray] at (7.4,3.1) {\((A,B,A)\)};
  % row labels
  \node[rlab] at (0.7,1.8) {\(I\)};
  \node[rlab] at (0.7,1.0) {\(J\)};
  \node[rlab] at (0.7,0.2) {\(K\)};
  % entries
  \node[cell] at (2.6,1.8) {B}; \node[cell] at (5.0,1.8) {B}; \node[cell] at (7.4,1.8) {A};
  \node[cell] at (2.6,1.0) {B}; \node[cell] at (5.0,1.0) {A}; \node[cell] at (7.4,1.0) {B};
  \node[cell] at (2.6,0.2) {B}; \node[cell] at (5.0,0.2) {A}; \node[cell] at (7.4,0.2) {A};
  % grid
  \draw[kwgray!60] (1.4,2.3) -- (8.6,2.3);
  \draw[kwgray!60] (1.4,1.4) -- (8.6,1.4);
  \draw[kwgray!60] (1.4,0.6) -- (8.6,0.6);
  \draw[kwgray!60] (1.4,-0.2) -- (8.6,-0.2);
  \draw[kwgray!60] (1.4,2.3) -- (1.4,-0.2);
  \draw[kwgray!60] (3.8,2.3) -- (3.8,-0.2);
  \draw[kwgray!60] (6.2,2.3) -- (6.2,-0.2);
  \draw[kwgray!60] (8.6,2.3) -- (8.6,-0.2);
  % S-duality arrows
  \draw[-{Stealth[length=2mm]}, thick, draw=kwred!70] (2.6,3.55) to[bend left=18] node[above=1pt, font=\scriptsize] {S-对偶} (5.0,3.55);
  \draw[-{Stealth[length=2mm]}, thick, draw=kwred!70] (5.0,3.9) to[bend left=18] node[above=1pt, font=\scriptsize] {S-对偶} (2.6,3.9);
  \node[font=\scriptsize, text=kwred!80!black, align=center] at (7.4,3.75) {S-对偶保持 \((A,B,A)\) 与 \((A,A,B)\)};
  % annotation
  \node[font=\footnotesize, text=kwgray, align=left, anchor=west] at (9.0,1.4) {脚注 24：不存在\\ \((B,B,A)\) 型膜；\\ 也不存在\\ \((A,A,A)\) 型膜。};
\end{tikzpicture}
\caption{三类关键膜相对于复结构 \(I,J,K\) 的 A/B 类型。每一行是一个复结构下的类型；每一列是一张膜。S-对偶交换 \((B,B,B)\) 与 \((B,A,A)\)，保持 \((A,B,A)\)（以及 \((A,A,B)\)）。这张表说明「\((A,B,A)\)」是相对于三个不同复结构的记录，而不是说这张膜在某个固定复结构下既 A 又 B。}
\label{fig:types}
\end{figure}

原文对「为什么只有这些组合」给的理由很短，我们把它补全（\pB）：脚注 24 断言不存在 \((B,B,A)\) 型膜——一个在复结构 \(I\) 与 \(J\) 中都全纯的子簇自动在 \(K\) 中全纯——也不存在 \((A,A,A)\) 型膜。\pA\anchor{src p.79} 我们在 \S\ref{subsec:cc} 已在 (11.4) 这一自然族内验证了 \((A,A,A)\) 不可能；脚注 24 的完整陈述是它的推广。\pC\anchor{eq 11.4} 于是四种可能的组合是 \((B,B,B)\)、\((B,A,A)\)、\((A,B,A)\)、\((A,A,B)\)，正好对应三类关键膜加上它的 K-复结构类比。\pA\anchor{src p.208}

\paragraph{排除错误直觉 7：第 12 节不是「再发明一遍膜」，也不产生新结论。}第 12 节的原文目标是「说明重要的膜可以在底层规范理论中定义」，并明确说这种提升不唯一：有些情形会给出多个不同的边界条件，它们在降到二维后（至少在奇点之外）看起来等价。\pA\anchor{src p.208} 所以第 12 节是验证，不是发现：它验证第 11 节的翻译所依赖的膜确实存在且有四维来源。把第 12 节读成「另一套独立构造」，会错过它与第 11 节之间的闭环关系。

\paragraph{排除错误直觉 8：「c.c. 膜是空间充满的，所以它的边界条件就是自由边界条件」。}空间充满只意味着支撑是整空间，不意味着边界条件自由。c.c. 膜的边界条件是混合型（12.31）--（12.32）：一部分方向 Dirichlet、一部分 Neumann，而混合的「斜切」由 target 规范场 \(f=\omega_J\)（12.34）决定。正是这个非零的 \(f\) 使得它的开弦代数变成微分算子层，而不是普通的函数层。\pA\anchor{src p.220}

\subsection{全章收束：翻译器如何闭合论证（S8 的产出）}
\label{subsec:close}

\paragraph{把链条接起来。}现在可以回答第 \ref{subsec:gap} 小节提出的缺口。全文的论证链是：

\begin{enumerate}
\item 第 8 节（S6）把 Wilson 线算子的作用翻译成「膜上的 Chan-Paton 层被张量积」，从而把 \(\pi_1(C)\to\LdualC\) 的同态等同于 electric eigenbrane；其 S-对偶像是卷在 Hitchin 纤维上的 magnetic eigenbrane。\pA\anchor{src p.185}
\item 第 9--10 节（S7）把 't Hooft 线算子等同于几何 Hecke 算子。\pA\anchor{src p.185}
\item 第 11 节（S8 的主体）造出 c.c. 膜，并证明它的开弦代数层是 \(\Dmod_{K_M^{1/2}}\)；于是「\(K\) 中的 A-膜」自动成为该代数层的模，翻译完成；\S11.3 再把 \(\Psi\) 推广，得到扭化微分算子，并在 \(\Psi=h\) 处退化为普通 D-模。\pA\anchor{src p.197}\pA\anchor{src p.207}
\item 第 12 节（S8 的检查）用四维规范理论的边界条件重新造出 \((B,B,B)\)、\((B,A,A)\)、\((A,B,A)\) 三类膜，并证明第 12.4 节的 \((A,B,A)\) 膜就是 c.c. 膜。\pA\anchor{src p.208}\pA\anchor{src p.221}
\end{enumerate}

拼起来：magnetic eigenbrane（\(K\) 中 A-膜）\(\mapsto\) \(\MC\) 上 D-模；而 Hecke 作用在膜侧由 't Hooft 算子给出，在 D-模侧由 Hecke 算子给出，于是得到了几何 Langlands 的标准表述——\(\MC\) 上同时是 Hecke 特征层的 D-模，对应 \(\pi_1(C)\to\LdualC\) 的局部系统。这就是本章兑现的 miracle M5：\emph{\(\MH\) 上存在典范的余各向同性 A-膜，使 A-膜范畴与 \(\MC\) 上的 D-模范畴在「模」的层面上联系起来}。

\paragraph{本章证明了什么，重建了什么。}为了避免读者高估或低估这一章，下面把结论按证据强度分类。

\begin{enumerate}
\item \textbf{A 级（原文明确证明）。}c.c. 膜的存在与唯一性（(11.1)--(11.5)、定义、单连通时唯一）；\((B_{c.c.},B_{c.c.})\) 弦层是 \(\Dmod_{K_M^{1/2}}\)（时间反演论证、整体截面）；翻译命题的框架（模结构、无瞬子修正、层化）；\(\Psi\) 推广与 \(f(\Psi)=\Psi\) 论证；第 12 节三类膜的边界条件与 c.c. 膜识别。\pA\anchor{src p.187}\pA\anchor{src p.195}\pA\anchor{src p.207}
\item \textbf{B 级（原文一句话带过、我们重建）。}(11.1) 的几何内容（原文引 [80]）；(11.4) 中 \(b\neq0\) 不可用的周期论证；正则量子化 (11.12)--(11.15)；转置算子给出 opposite 代数的分部积分；\(G=U(1)\) 中「协变导数给出平坦联络」；\((A,A,A)\) 不存在的自然族验证；Hecke 作用与翻译的交换性。\pB\anchor{src p.186}\pB\anchor{src p.194}\pB\anchor{src p.201}
\item \textbf{C 级（我们判断、原文未明言）。}(11.1) 的功能是「让翻译可积」而非仅「让膜存在」；非交换形变发生在余切方向而非 \(\MC\) 本身；\(\Psi\neq0\) 的几何变化不可视为装饰；第 12 节与第 11.2 节的对应是检查而非独立证明。\pC\anchor{src p.189}\pC\anchor{src p.221}
\end{enumerate}

\paragraph{原文的省略清单。}把散布在本章的「重建点」集中列出，便于读者核对（每一条都给出锚点，供对照原文页文本）：

\begin{enumerate}
\item (11.1) 的推导：原文只引 [80]，我们补出「\(\omega-\ii F\) 是 (2,0) 型全纯辛形式」的运算（\S\ref{subsec:cc}，\pB）。\anchor{src p.186}
\item (11.4) 的周期论证：原文一句「generally does not have periods which are integral multiples of \(2\pi\)」，我们补出「\(\omega_I\) 在 Hitchin 纤维上有非零周期」的算术（\S\ref{subsec:cc}，\pB）。\anchor{src p.187}
\item (11.12)--(11.15)：我们补出 Weyl 排序下的正则量子化（\S\ref{subsec:translation}，\pB）。\anchor{src p.192}
\item opposite 代数的分部积分：原文说「证法留给 [140]」，我们写出转置算子的式子（\S\ref{subsec:translation}，\pB）。\anchor{src p.194}
\item \(G=U(1)\) 的「平坦」：原作在写出 (11.24)--(11.25) 后直接说定义了平坦联络，我们补出「曲率由 \(\upsilon\) 的调和性保证为零」的一步（\S\ref{subsec:translation}，\pB）。\anchor{src p.201}
\item 第 12.4 节的边界可观测量与第 11.2 节环层的对应：原文在两处分别得到结论，没有显式对照；我们把它们列为检查点（\S\ref{subsec:section12}，\pB）。\anchor{src p.221}
\end{enumerate}

\paragraph{本章的局限与最薄弱的两处。}如实列出，供讨论班继续工作。

\begin{enumerate}
\item \textbf{Hecke 作用的函子性没有逐条证明。}第 11 节证明「膜 \(\mapsto\) D-模」，第 9--10 节证明「't Hooft \(\mapsto\) Hecke」，拼起来需要「翻译函子与 Hecke 作用交换」的图；原文只以「都在同一个四维 TFT 中」的自然性带过。我们判断这是最薄弱的一环：若交换性不成立，得到的 D-模未必是 Hecke 特征层。\pC\anchor{src p.185}
\item \textbf{非整数 \(\Psi\) 的代数层只由微扰论证确定。}原文的 \(f(\Psi)=\Psi\) 依赖把转移函数算到 \(1/(\Im\tau)^2\) 阶、并利用微扰展开线性于 \(\Re\tau\)，这不是非微扰构造；完全层化还有收敛问题（脚注 58 的代数结构选择）。\pB\anchor{src p.207}\pA\anchor{src p.206}
\end{enumerate}

此外还有原文自己点出的两处限制：其一，全文只处理非分歧（unramified）情形，没有讨论分歧的几何 Langlands；其二，某些陈述对 \(C\) 的 spin 结构敏感，原文在 \(\Psi=h\) 的普通 D-模陈述里明确提到要对某些 \(G\) 仔细考虑 [5,26]。\pA\anchor{src p.207}

\begin{keyidea}
\textbf{一句话总结本章。}c.c. 膜是全文的翻译器：因为它在 \(K\) 中是空间充满的 A-膜、且它的 \(F\) 由 \(\omega_J\) 典范地确定，它的开弦代数恰好是全纯微分算子层 \(\Dmod_{K_M^{1/2}}\)；于是「\(K\) 中 A-膜」自动成为「\(\MC\) 上 D-模」。第 11.3 节把 \(\Psi=0\) 推广到任意 \(\Psi\)，第 12 节用四维边界条件为三类膜提供独立来源；两者一起把第 8 节的「特征膜 \(\leftrightarrow\) 局部系统」升级为几何 Langlands 的标准表述「D-模 \(\leftrightarrow\) 局部系统」。
\end{keyidea}
